%% =====================================================================
%%  Capítulo 8 -- Principio de invarianza y construcción de funciones de Lyapunov
%%  Formato nuevo (noviembre de 2026). Integra: LaSalle/LaSalle_Invarianza.tex y
%%  Construccion/Construccion_Lyapunov.tex (paquetes de octubre de 2026 para
%%  Handout7, que acompañaban las láminas 97--108 de Lyapunov2026) y los scripts
%%  Codigos/LaSalle/lasalle_*.m y Codigos/GradienteZubov/construccion_*.m,
%%  reescritos en codigo/cap08_invarianza. Lo que ya está en el capítulo 7
%%  (Barbashin--Krasovskii, la trampa radial, el término cruzado del péndulo,
%%  la ecuación de Lyapunov con Q = C'C, la estimación cuadrática de van der Pol)
%%  se cita y no se repite. Correcciones: ver MIGRACION.md.
%%  Números verificados con herramientas/cap08_numeros.py y
%%  codigo/cap08_invarianza/verificarCap08.m.
%% =====================================================================

\capitulo[cap:invarianza]{Principio de invarianza y construcción de funciones de Lyapunov}
  {El tiempo es un río que me arrebata, pero yo soy el río.}
  {Jorge Luis Borges}

\begin{objetivos}[partes/08-invarianza.html]
  \item Distinguir conjuntos invariantes de positivamente invariantes, construir un conjunto compacto y positivamente invariante con una componente de un conjunto de nivel, y usar las propiedades del conjunto $\omega$-límite.
  \item Enunciar y probar el principio de invarianza de LaSalle, calcular $E$ y su mayor subconjunto invariante $M$ con la prueba de permanencia, y concluir exactamente lo que el teorema permite (y no más).
  \item Reconocer qué puede ser $M$ (un punto, un ciclo límite, una recta de equilibrios, un subespacio no observable) y por qué la convergencia a $M$ no da por sí sola convergencia a un punto ni estabilidad.
  \item Estimar la región de atracción con conjuntos de nivel cuando $\Vdot$ es solo semidefinida.
  \item Construir funciones de Lyapunov con el gradiente variable y con el método de Krasovskii, y caracterizar la región de atracción exacta con la ecuación de Zubov.
\end{objetivos}

%% ---------------------------------------------------------------------
\section{Lo que el capítulo~7 dejó sin resolver}
\label{sec:iv-porque}
%% ---------------------------------------------------------------------

El \cref{cap:lyapunov} terminó con dos preguntas abiertas. La primera apareció con el péndulo con fricción (\cref{ej:ly-pendenergia}): la energía da $\Vdot=-b\,x_2^2\le0$, que se anula en todo el eje $x_2=0$, y el método directo solo concluye \emph{estabilidad}, aunque la física y la simulación dicen que el péndulo termina en reposo. La segunda apareció al estimar la región de atracción de van der Pol en tiempo invertido (\cref{ej:ly-vdp}): una función cuadrática no puede seguir la forma del ciclo límite, y la mejor elipse cubre menos de la mitad de la región. Este capítulo responde las dos:
\begin{enumerate}
  \item Si $\Vdot\le0$ y $\Vdot$ se anula en algo más que el origen, ¿qué se puede concluir, y con qué argumento? $\to$ \clave{principio de invarianza de LaSalle}, \S\ref{sec:iv-invariantes}--\ref{sec:iv-region}.
  \item ¿De dónde sale una $V$ cuando no hay una energía evidente, o cuando la energía se queda corta? $\to$ \clave{gradiente variable}, \clave{Krasovskii} y \clave{Zubov}, \S\ref{sec:iv-gradiente}--\ref{sec:iv-zubov}.
\end{enumerate}

Conviene ver con precisión qué da, y qué no da, la desigualdad $\Vdot\le0$. Para el péndulo
\begin{equation}
  \dot x_1=x_2,\qquad \dot x_2=-a\sin x_1-b\,x_2,\qquad V(x)=a(1-\cos x_1)+\tfrac12x_2^2,\qquad \Vdot=-b\,x_2^2,
  \label{eq:iv-pendulo}
\end{equation}
que es \eqref{eq:ly-pendulo} con $a=g/l$ y $b=k/m$, se obtienen, con rigor, tres hechos:
\begin{enumerate}
  \item $t\mapsto V(x(t))$ es no creciente y acotada por debajo por $0$: existe $V_\infty=\lim_{t\to\infty}V(x(t))$.
  \item Si $x_0$ está en un conjunto de nivel acotado de $V$ (en la \cref{sec:iv-invariantes} se precisa cuál), la solución no sale de él: está acotada.
  \item Integrando, $b\int_0^\infty x_2(\tau)^2\dif\tau=V(x_0)-V_\infty<\infty$.
\end{enumerate}
Ninguno dice que $x(t)\to0$. El tercero sugiere que $x_2$ es pequeña ``casi todo el tiempo'', pero una función de cuadrado integrable no tiene por qué tender a cero: pulsos de altura $1$ y ancho $1/k^2$ en $t=k$ tienen $\int x_2^2=\sum_k1/k^2<\infty$ y no se apagan (en el \cref{cap:avanzada} se vuelve sobre esto, \cref{tr:av-integral}); y aun si $x_2\to0$, falta $x_1$. La pieza que falta es un argumento sobre las soluciones que pueden \emph{quedarse} donde $\Vdot=0$: para que $V$ dejara de decrecer, la trayectoria tendría que permanecer en el eje $x_2=0$; si $x_2(t)\equiv0$, entonces $\dot x_2(t)\equiv0$, y la segunda ecuación da $\sin x_1(t)\equiv0$, es decir $x_1(t)\equiv0$ si $\abs{x_1}<\pi$. El único movimiento que cabe en $\{\Vdot=0\}$ es el reposo. El argumento es correcto, pero deja tres preguntas que el teorema de LaSalle responde con precisión: ¿por qué la trayectoria tiene que ``terminar'' en algún lado? ¿Por qué basta mirar las soluciones que \emph{permanecen} en $\{\Vdot=0\}$? ¿En qué región vale todo esto?

\begin{figure}[ht]
\centering
\includegraphics[width=0.95\linewidth]{cap08/t_mapa_invarianza}
\caption{Mapa del capítulo. El principio de invarianza cierra el hueco de $\Vdot\le0$; los métodos de construcción dan funciones (a veces estrictas) donde no hay una energía evidente. Las dos líneas se encuentran al estimar la región de atracción.}
\label{fig:iv-mapa}
\end{figure}

\notamargen{\textbf{Ruta de las clases.} Semana~10, después de la región de atracción del \cref{cap:lyapunov}: \S\ref{sec:iv-porque}--\ref{sec:iv-ejemplos}. Semana~11: \S\ref{sec:iv-region}--\ref{sec:iv-zubov}. Las pruebas del \cref{lem:iv-subnivel} y del \cref{thm:iv-lasalle} se hacen en clase.}

\notamargen{\textbf{Dos usos de ``invariante''.} El curso dice que un sistema es ``invariante en el tiempo'' cuando $f$ no depende de $t$: $\dot x=f(x)$. Un \emph{conjunto invariante} es otra cosa: uno del que las soluciones no salen. El principio de invarianza habla de lo segundo y vale para sistemas invariantes en el tiempo.}

\begin{historia}[Barbashin, Krasovskii, LaSalle]
Los dos corolarios que más se usan (estabilidad asintótica, local y global, con $\Vdot\le0$) se deben a E.~A.~Barbashin y N.~N.~Krasovskii en 1952~\citep{BaKr52} y aparecen en la monografía de Krasovskii de 1959~\citep{Kra63}. J.~P.~LaSalle formuló en 1960 la versión general en términos de conjuntos invariantes~\citep{LaS60}. Por eso se encuentran los nombres ``principio de LaSalle'', ``teorema de Krasovskii--LaSalle''~\citep{KeBr23} y ``teorema de Barbashin y Krasovskii'' para los corolarios~\citep{Kha02}. Los métodos de construcción son de la misma época: el gradiente variable, de Schultz y Gibson en 1962~\citep{ScGi62}; el método de Krasovskii, de su monografía; la ecuación de Zubov, de 1957~\citep{Zub64}. Hoy la construcción se hace también con optimización convexa, buscando $V$ como suma de cuadrados de polinomios~\citep{Par00}.
\end{historia}

\begin{pausa}
Antes de seguir, prediga: (1) si $V$ es definida positiva, $\Vdot\le0$ y la solución está acotada, ¿converge la solución a algún punto donde $\Vdot=0$? (2) Si $\Vdot\le0$ y los únicos equilibrios en $\{\Vdot=0\}$ están en el origen, ¿convergen las soluciones al origen? (3) Si todas las soluciones de un anillo convergen a un equilibrio, ¿es estable ese equilibrio? (4) ¿Puede una función que \emph{no} es definida positiva probar que algo converge?
\respuesta{(1) No necesariamente a un \emph{punto}: converge a un \emph{conjunto} invariante donde $\Vdot=0$, que puede ser un ciclo límite o una recta de equilibrios (\cref{ej:iv-ciclo,ej:iv-adaptativo}). (2) No: $\{\Vdot=0\}$ puede contener órbitas periódicas sin ningún equilibrio, y entonces la solución puede quedarse oscilando (\cref{tr:iv-equilibrios}). (3) No: converger no es ser estable (\cref{ej:iv-inestable}). (4) Sí: el \cref{thm:iv-lasalle} no pide que $V$ sea definida positiva, y con $V=\tfrac14(r^2-1)^2$ prueba que un ciclo límite atrae (\cref{ej:iv-ciclo}).}
\end{pausa}

%% ---------------------------------------------------------------------
\section{Conjuntos invariantes y conjuntos límite}
\label{sec:iv-invariantes}
%% ---------------------------------------------------------------------

En todo el capítulo se considera el sistema invariante en el tiempo
\begin{equation}
  \dot x=f(x),\qquad f:\D\to\R^n\ \text{localmente Lipschitz},
  \label{eq:iv-sistema}
\end{equation}
con $\D\subset\R^n$ un dominio (abierto y conexo), y se escribe $\flujo t{x_0}$ para la solución que en $t=0$ pasa por $x_0$, como en el \cref{cap:orbitas}. Para $V:\D\to\R$ de clase $\Cont1$, $\Vdot(x)=\jac Vx(x)\,f(x)$.

\subsection{Invariante y positivamente invariante}

\begin{definicion}[label=def:iv-invariante]{Conjuntos invariantes \citep[\S4.2]{Kha02}}
Un conjunto $M\subset\D$ es \clave{invariante} respecto a \eqref{eq:iv-sistema} si
\[
  x_0\in M\implica\flujo t{x_0}\in M\qquad\text{para todo }t\in\R ,
\]
y es \clave{positivamente invariante} si la implicación vale para todo $t\ge0$.
\end{definicion}

La diferencia importa: \alerta{invariante quiere decir que el conjunto está hecho de trayectorias completas, hacia el futuro y hacia el pasado}. Un equilibrio, una órbita periódica y una órbita homoclínica con su silla son invariantes, y la unión de conjuntos invariantes es invariante. En cambio, un conjunto de nivel $\{V\le c\}$ con $\Vdot\le0$ es, en general, solo positivamente invariante: hacia atrás en el tiempo la energía del péndulo crece y la solución sale del conjunto. Ya se vio un ejemplo en la recta: en la dinámica halcón--paloma, el intervalo $[0,1]$ es invariante porque sus extremos son equilibrios (\cref{sec:un-halcon}).

\begin{lema}[label=lem:iv-subnivel]{Componentes de un conjunto de nivel}
Sea $V:\D\to\R$ de clase $\Cont1$ con $\Vdot(x)\le0$ en $\D$. Si $\Nivel c$ es una componente conexa de $\{x\in\D:\ V(x)\le c\}$, acotada y con $\overline{\Nivel c}\subset\D$, entonces $\Nivel c$ es compacta y positivamente invariante.
\end{lema}

\begin{demostracion}
Sea $x_0\in\Nivel c$ y $[0,t^*)$ el intervalo maximal en el que $\flujo t{x_0}\in\D$. Allí $V(\flujo t{x_0})\le V(x_0)\le c$, así que la solución permanece en $\{V\le c\}$; como $t\mapsto\flujo t{x_0}$ es continua y $\flujo0{x_0}\in\Nivel c$, permanece en la \emph{misma} componente conexa (una curva continua no salta de una componente a otra). Como $\Nivel c$ es acotada y $\overline{\Nivel c}\subset\D$, la solución no puede salir de $\D$ ni escapar en tiempo finito (\cref{cap:preliminares}): $t^*=\infty$. Por último, $\Nivel c$ es cerrada en $\D$ y su clausura está en $\D$, luego es cerrada en $\R^n$; con la acotación, es compacta.
\end{demostracion}

\begin{trampa}[label=tr:iv-subnivel]{``$\{V\le c\}$ es acotado y todo punto con $V(x_0)\le c$ está en $\Nivel c$''}
\plausible{$V$ es definida positiva cerca del origen, así que sus conjuntos de nivel pequeños son ``bolitas'' alrededor del origen.}
Para el péndulo \eqref{eq:iv-pendulo}, $\{x\in\R^2:\ V(x)\le c\}$ \emph{nunca} es acotado, por pequeño que sea $c>0$: como $V$ es periódica en $x_1$, contiene una isla alrededor de cada $(2k\pi,0)$ (zonas grises de la \cref{fig:iv-pendulo}). Un punto con $V(x_0)<c$ en otra isla no está en $\Nivel c$, y en el \cref{ej:iv-duffing} un punto así escapa a infinito.
\correcto{Se toma la \emph{componente conexa} de $\{V\le c\}$ que contiene al origen (para el péndulo con $c<2a$, la que está en $\abs{x_1}<\pi$). \citet[\S4.1]{Kha02} evita el problema con $\Omega_\beta=\{x\in\Bola r:\ V(x)\le\beta\}$ y $\beta<\min_{\norm x=r}V(x)$.}
\end{trampa}

\subsection{Conjuntos límite}

El conjunto $\omega$-límite se definió en el \cref{cap:orbitas} (\cref{def:oc-omega}): $p\in\omegalim{x_0}$ si existe $t_k\to\infty$ con $\flujo{t_k}{x_0}\to p$. Para LaSalle hace falta, además, el lenguaje de convergencia a un conjunto: $\flujo t{x_0}$ \clave{tiende a} $M$ si
\[
  \lim_{t\to\infty}\dist\bigl(\flujo t{x_0},M\bigr)=0,\qquad \dist(p,M)=\inf_{q\in M}\norm{p-q} .
\]
El \cref{lem:oc-omega} dice que si la solución es \emph{acotada} y se queda en $\D$ con su clausura, $\omegalim{x_0}$ es no vacío, compacto, conexo e \alerta{invariante}, y la solución tiende a él. La invarianza es la propiedad que hace funcionar el teorema de LaSalle. Vale la pena ver de dónde sale: si $\flujo{t_k}{x_0}\to p$ y $\tau$ es fijo, la dependencia continua de las soluciones respecto a la condición inicial en el intervalo compacto $[0,\tau]$ (\cref{cap:preliminares}) da $\flujo{t_k+\tau}{x_0}=\flujo\tau{\flujo{t_k}{x_0}}\to\flujo\tau p$, así que $\flujo\tau p\in\omegalim{x_0}$. La solución que pasa por un punto límite se queda entre los puntos límite.

\begin{ejemplo}[label=ej:iv-limite]{Tres conjuntos límite}
Encuentre $\omegalim{x_0}$ para (a) el oscilador amortiguado $\dot x_1=x_2$, $\dot x_2=-x_1-0.4x_2$; (b) van der Pol con $\mu=1$ desde $x_0=(0.1,0)$; (c) $\dot x_1=0.5$, $\dot x_2=-x_2$.
\solucion (a) La solución converge al origen: $\omegalim{x_0}=\{0\}$. (b) La solución se acerca al ciclo límite $\Gamma$ (\cref{cap:orbitas}), de amplitud $2.009$ en $x_1$, pero $x_1(t)$ sigue oscilando entre $\pm2.009$: $\lim_{t\to\infty}\flujo t{x_0}$ no existe, y sin embargo $\omegalim{x_0}=\Gamma$. Las muestras $\flujo n{x_0}$, $n=1,2,\dots$, se acumulan sobre todo el ciclo (\cref{fig:iv-limite}, centro). (c) $x_2\to0$, pero $x_1=x_1(0)+0.5t\to\infty$: no hay ninguna sucesión convergente y $\omegalim{x_0}=\emptyset$.
\end{ejemplo}

\begin{figure}[ht]
\centering
\includegraphics[width=\linewidth]{cap08/f_limite}
\caption{Tres conjuntos $\omega$-límite. (a) Un punto. (b) Van der Pol: las muestras $\flujo n{x_0}$ (color según $n$) se acumulan sobre el ciclo (coral), y la solución tiende al \emph{conjunto} sin tender a ningún punto. (c) Sin acotación no hay conjunto límite.}
\label{fig:iv-limite}
\end{figure}

\codigorepo{cap08_invarianza/C8_Ej1_ConjuntosLimite.m}{limite}{Muestras $\flujo n{x_0}$ y los tres conjuntos límite}

El panel (b) explica por qué LaSalle habla de convergencia a un \emph{conjunto}; el panel (c), por qué la compacidad es indispensable.

\begin{pausa}
Para $\dot r=r(1-r^2)$, $\dot\theta=1$, ¿cuál es $\omegalim{x_0}$ si $x_0=0$? ¿Y si $x_0\neq0$? ¿Cuál es el $\alpha$-límite de un $x_0$ con $0<\norm{x_0}<1$?
\respuesta{$\omegalim0=\{0\}$ (es un equilibrio). Si $x_0\neq0$, $r(t)\to1$ y $\omegalim{x_0}$ es el círculo unitario. Hacia atrás, si $0<r_0<1$, $r(t)\to0$ cuando $t\to-\infty$: $\alphalim{x_0}=\{0\}$.}
\end{pausa}

%% ---------------------------------------------------------------------
\section{El principio de invarianza}
\label{sec:iv-lasalle}
%% ---------------------------------------------------------------------

\begin{ideaclave}
\textbf{La idea en una frase.} Si $V$ no puede crecer a lo largo de las soluciones y estas permanecen en un compacto, cada solución se acumula en un conjunto donde $V$ es \emph{constante}. Ese conjunto es invariante, así que está contenido en $\{\Vdot=0\}$ \emph{y está hecho de soluciones completas}. Si la única solución completa que cabe en $\{\Vdot=0\}$ es el equilibrio, todas convergen a él.
\end{ideaclave}

\subsection{El conjunto $E$ y su mayor subconjunto invariante $M$}

\begin{definicion}[label=def:iv-EM]{Los conjuntos $E$ y $M$}
Sean $\Omega\subset\D$ y $V$ de clase $\Cont1$ con $\Vdot\le0$ en $\Omega$. Se definen
\[
  E=\{x\in\Omega:\ \Vdot(x)=0\}
\]
y $M$, el \clave{mayor conjunto invariante contenido en $E$}: la unión de todos los subconjuntos invariantes de $E$. Equivalentemente,
\[
  x_0\in M\iff\flujo t{x_0}\in E\quad\text{para todo }t\in\R .
\]
\end{definicion}

$M$ está bien definido porque la unión de conjuntos invariantes es invariante (puede ser vacío). Siempre $M\subset E$, y en general la inclusión es estricta: en el péndulo, $E$ es un segmento del eje $x_2=0$ y $M=\{0\}$.

\paragraph{Cómo se calcula $M$: la prueba de permanencia.} Suponga que $E=\{x\in\Omega:\ h(x)=0\}$ (en el péndulo, $h(x)=x_2$). Si una solución permanece en $E$, entonces $h(x(t))\equiv0$, y por lo tanto \emph{todas} sus derivadas se anulan:
\begin{equation}
  h(x(t))\equiv0\implica\Lie fh\bigl(x(t)\bigr)\equiv0\implica\Lie f^2h\bigl(x(t)\bigr)\equiv0\implica\cdots,
  \label{eq:iv-permanencia}
\end{equation}
con $\Lie fh=\jac hx f$. Cada derivada impone una restricción algebraica nueva; se deriva hasta que las restricciones obliguen a $x(t)$ a estar en un conjunto que ya se reconoce (típicamente, un equilibrio).\notamargen{En el \cref{cap:linealizacion} las $\Lie f^kh$ se llamarán derivadas de Lie y se usarán en la linealización por retroalimentación. Calcular $M$ es preguntar qué estados son \emph{inobservables} desde la ``salida'' $y=h(x)$: \cref{ej:iv-observable}.} Para el péndulo:
\[
  x_2\equiv0\implica\dot x_2=-a\sin x_1-b\,x_2\equiv0\implica\sin x_1\equiv0\implica x_1\equiv0\qquad(\abs{x_1}<\pi).
\]
La geometría es la de la \cref{fig:iv-pendulo}: sobre $E$ el campo vale $f(x_1,0)=(0,-a\sin x_1)$, es vertical y, salvo en el origen, \emph{transversal} a $E$. Una solución que toca $E$ lo atraviesa de inmediato.

\begin{trampa}[label=tr:iv-EM]{``Como $\Vdot=0$ en el eje $x_2=0$, la solución termina en algún punto del eje''}
\plausible{La solución se acumula donde $\Vdot=0$, y $\Vdot=0$ en todo el eje.}
Termina en $M$, no en $E$. Los puntos de $E\setminus M$ son puntos por los que se \emph{pasa}, no puntos donde se \emph{queda}: la trayectoria de la \cref{fig:iv-pendulo} cruza el eje doce veces en $[0,40]$ y en cada cruce $V$ sigue bajando.
\correcto{Hay que calcular $M$ con la prueba de permanencia \eqref{eq:iv-permanencia}.}
\end{trampa}

\begin{trampa}[label=tr:iv-equilibrios]{``Los únicos equilibrios en $E$ están en el origen, luego $M=\{0\}$''}
\plausible{Una solución que se queda quieta en $E$ es un equilibrio; si no hay más equilibrios en $E$, solo queda el origen.}
Una solución puede quedarse en $E$ \emph{sin estar quieta}: $E$ puede contener órbitas periódicas, o subespacios llenos de ellas (\cref{ej:iv-ciclo,ej:iv-observable}). Con este razonamiento, en el \cref{ej:iv-ciclo} se concluiría $M=\emptyset$.
\correcto{La pregunta es qué \emph{soluciones} pueden permanecer en $E$, no qué \emph{puntos} de $E$ son estacionarios.}
\end{trampa}

\subsection{El teorema}

\begin{teorema}[label=thm:iv-lasalle]{Principio de invarianza de LaSalle \citep{LaS60}, \citep[Teo.~4.4]{Kha02}}
Considere \eqref{eq:iv-sistema}. Sea $\Omega\subset\D$ un conjunto \clave{compacto} y \clave{positivamente invariante}, y $V:\D\to\R$ de clase $\Cont1$ con $\Vdot(x)\le0$ en $\Omega$. Sea $E=\{x\in\Omega:\ \Vdot(x)=0\}$ y $M$ el mayor conjunto invariante contenido en $E$. Entonces toda solución que comienza en $\Omega$ tiende a $M$ cuando $t\to\infty$.
\end{teorema}

\begin{demostracion}
Sea $x_0\in\Omega$ y $x(t)=\flujo t{x_0}$. (i) \emph{La solución existe para todo $t\ge0$ y se queda en $\Omega$}: por la invarianza positiva, $x(t)\in\Omega$ mientras exista, y como $\Omega$ es compacto no hay escape en tiempo finito. (ii) \emph{$V(x(t))$ tiene límite}: es no creciente porque $\Vdot\le0$ en $\Omega$, y acotada por debajo porque $V$ es continua en el compacto $\Omega$; sea $\ell=\lim_{t\to\infty}V(x(t))$. (iii) \emph{El conjunto límite existe y está en $\Omega$}: por el \cref{lem:oc-omega}, $\omegalim{x_0}$ es no vacío, compacto e invariante, y $x(t)\to\omegalim{x_0}$; como $\Omega$ es cerrado, $\omegalim{x_0}\subset\Omega$. (iv) \emph{$V\equiv\ell$ en $\omegalim{x_0}$, y por eso $\omegalim{x_0}\subset M$}: si $x(t_k)\to p$, por continuidad $V(p)=\ell$. Como $\omegalim{x_0}$ es invariante, la solución que pasa por $p$ se queda en $\omegalim{x_0}$, donde $V\equiv\ell$; luego $\Vdot(p)=\frac{\dif}{\dif t}V(\flujo tp)\big|_{t=0}=0$. Así $\omegalim{x_0}\subset E$, y como es invariante, $\omegalim{x_0}\subset M$. (v) \emph{Conclusión}: $\dist(x(t),M)\le\dist(x(t),\omegalim{x_0})\to0$.
\end{demostracion}

La prueba entera cabe en una cadena de inclusiones, y cada eslabón usa una hipótesis distinta:
\begin{equation}
  \omegalim{x_0}\ \subset\ M\ \subset\ E\ \subset\ \Omega
  \qquad\begin{array}{l}
  \text{compacidad}\implica\omegalim{x_0}\neq\emptyset,\\
  V\text{ constante en }\omegalim{x_0}\implica\omegalim{x_0}\subset E,\\
  \text{invarianza de }\omegalim{x_0}\implica\omegalim{x_0}\subset M .
  \end{array}
  \label{eq:iv-cadena}
\end{equation}

Conviene leer con cuidado lo que el teorema \emph{no} pide y lo que \emph{no} concluye.
\begin{itemize}
  \item \alerta{No pide que $V$ sea definida positiva}, ni que $\Omega$ sea un conjunto de nivel de $V$: basta que $V$ sea $\Cont1$ y no crezca en $\Omega$. Esto permite probar la atracción de un ciclo límite (\cref{ej:iv-ciclo}).
  \item \alerta{Sí pide un $\Omega$ compacto y positivamente invariante}, y su existencia hay que establecerla aparte: con el \cref{lem:iv-subnivel}, con una región atrapante o de otra forma. Sin compacidad falla el paso (iii).
  \item \alerta{Concluye convergencia a un conjunto, no a un punto, y no concluye estabilidad de nada.} Las dos cosas requieren argumentos adicionales (\cref{ej:iv-adaptativo,ej:iv-inestable}).
  \item \alerta{Es un resultado para sistemas invariantes en el tiempo.} El paso (iv) usa que $\omegalim{x_0}$ es invariante, lo que en general falla para $\dot x=f(t,x)$; allí se usan Barbalat (\cref{lem:vt-barbalat}) o LaSalle--Yoshizawa (\cref{thm:vt-ly}).
\end{itemize}

\subsection{Corolarios: estabilidad asintótica local y global}

Cuando lo que se quiere es la estabilidad asintótica del origen, el teorema se especializa en los resultados de Barbashin y Krasovskii.

\begin{corolario}[label=cor:iv-as]{Estabilidad asintótica con $\Vdot\le0$ \citep[Cor.~4.1]{Kha02}}
Sea $x=0$ un equilibrio de \eqref{eq:iv-sistema} y $V:\D\to\R$ de clase $\Cont1$, \clave{definida positiva} en un dominio $\D$ que contiene al origen, con $\Vdot\le0$ en $\D$. Sea $S=\{x\in\D:\ \Vdot(x)=0\}$ y suponga que \clave{ninguna solución distinta de $x(t)\equiv0$ puede permanecer idénticamente en $S$}. Entonces el origen es asintóticamente estable.
\end{corolario}

\begin{demostracion}
\emph{Estabilidad}: es el \cref{thm:ly-directo}. \emph{Atracción}: tome $r>0$ con $\Bola r\subset\D$, $\alpha=\min_{\norm x=r}V(x)>0$, $\beta\in(0,\alpha)$ y $\Omega_\beta=\{x\in\Bola r:\ V(x)\le\beta\}$. Una solución que arranca en $\Omega_\beta$ no puede alcanzar la esfera $\norm x=r$, donde $V\ge\alpha>\beta$, porque $V(x(t))$ no crece: $\Omega_\beta$ es compacto y positivamente invariante, y contiene una bola $\Bola\delta$ por la continuidad de $V$ en el origen. Por el \cref{thm:iv-lasalle}, toda solución que arranca en $\Omega_\beta$ tiende a $M\subset S\cap\Omega_\beta$, y por hipótesis $M=\{0\}$.
\end{demostracion}

\begin{corolario}[label=cor:iv-gas]{Estabilidad asintótica global con $\Vdot\le0$ \citep[Cor.~4.2]{Kha02}}
Sea $x=0$ un equilibrio de \eqref{eq:iv-sistema} con $\D=\R^n$, y $V:\R^n\to\R$ de clase $\Cont1$, definida positiva y \clave{radialmente no acotada}, con $\Vdot\le0$ en $\R^n$. Si ninguna solución distinta de $x(t)\equiv0$ puede permanecer idénticamente en $S=\{x:\ \Vdot(x)=0\}$, el origen es globalmente asintóticamente estable.
\end{corolario}

\begin{demostracion}
Dado $x_0$, sea $c=V(x_0)$. Por ser $V$ radialmente no acotada, $\Omega=\{x\in\R^n:\ V(x)\le c\}$ es acotado; es cerrado y, como $\Vdot\le0$ en todo $\R^n$, positivamente invariante. Se aplica el \cref{thm:iv-lasalle} con ese $\Omega$, que no necesita ser conexo.
\end{demostracion}

Si $\Vdot$ es definida negativa, $S=\{0\}$ y los corolarios son los \cref{thm:ly-directo,thm:ly-bk}. La hipótesis radial no sobra: el sistema de la \cref{tr:ly-radial} tiene $V$ definida positiva, $\Vdot<0$ fuera del origen (así que $S=\{0\}$) y aun así hay soluciones que escapan. Sin $\Omega$ compacto no hay conjunto límite garantizado: para esas soluciones, $\omegalim{x_0}=\emptyset$.

\begin{trampa}[label=tr:iv-pasar]{``Las soluciones pasan por $S$, así que la hipótesis del corolario falla''}
\plausible{El corolario pide que las soluciones no se queden en $S$, y en el péndulo todas llegan a $S$ una y otra vez.}
La hipótesis no dice que las soluciones eviten $S$; dice que ninguna solución no trivial puede quedarse en $S$ \emph{para todo tiempo}. Las del péndulo cruzan $S$ infinitas veces y el corolario se aplica.
\correcto{``Pasar por'' no es ``permanecer en''. Se verifica con la prueba de permanencia.}
\end{trampa}

\begin{receta}[Cómo aplicar LaSalle]
\begin{enumerate}
  \item \textbf{Escoger $V$ y calcular $\Vdot$.} Verificar $\Vdot\le0$ en la región de interés. ($V$ no tiene que ser definida positiva para el teorema; sí para los corolarios.)
  \item \textbf{Exhibir $\Omega$ compacto y positivamente invariante} que contenga la condición inicial: la componente conexa de $\{V\le c\}$ que la contiene (\cref{lem:iv-subnivel}), $\Omega_\beta=\{x\in\Bola r:\ V\le\beta\}$, $\{V\le V(x_0)\}$ si $V$ es radialmente no acotada, o una región atrapante obtenida de otra forma.
  \item \textbf{Calcular $E=\{x\in\Omega:\ \Vdot(x)=0\}$.}
  \item \textbf{Calcular $M$ con la prueba de permanencia} \eqref{eq:iv-permanencia}: suponer $x(t)\in E$ para todo $t$ y derivar las restricciones que definen $E$.
  \item \textbf{Concluir con precisión.} Siempre: $x(t)\to M$ para $x_0\in\Omega$. Si $M=\{0\}$ y $V$ es definida positiva: estabilidad asintótica, con $\Omega$ dentro de la región de atracción. Si además $V$ es radialmente no acotada y $\Vdot\le0$ en $\R^n$: estabilidad asintótica global. Si $M$ es un conjunto de equilibrios u órbitas: solo convergencia a ese conjunto.
\end{enumerate}
\end{receta}

\begin{ejemplo}[label=ej:iv-pendulo]{El péndulo con fricción, por fin}
Pruebe que el origen de \eqref{eq:iv-pendulo} con $a,b>0$ es asintóticamente estable y dé un estimado de su región de atracción.
\solucion \emph{Paso 1.} En $\D=(-\pi,\pi)\times\R$, $V$ es definida positiva y $\Vdot=a\sin x_1\,x_2+x_2(-a\sin x_1-b\,x_2)=-b\,x_2^2\le0$.
\emph{Paso 2.} Para $c<2a$ (la energía del péndulo invertido), en las rectas $x_1=\pm\pi$ se tiene $V\ge2a>c$, así que la componente $\Nivel c$ de $\{V\le c\}$ que contiene al origen no las cruza y está en $\abs{x_1}<\pi$; como además $x_2^2\le2c$, es acotada, está en $\D$ y, por el \cref{lem:iv-subnivel}, es compacta y positivamente invariante.
\emph{Paso 3.} $E=\{(x_1,0):\ \abs{x_1}\le\arccos(1-c/a)\}$; con $a=1$ y $c=1.8$, $\abs{x_1}\le2.498$.
\emph{Paso 4.} $x_2\equiv0$ obliga $a\sin x_1\equiv0$ y, como $\abs{x_1}<\pi$, $x_1\equiv0$: $M=\{0\}$.
\emph{Paso 5.} Por el \cref{cor:iv-as}, el origen es asintóticamente estable, y todo $\Nivel c$ con $c<2a$ está en la región de atracción: toda condición inicial con $\abs{x_1(0)}<\pi$ y energía menor que $2a$ termina en reposo. Con $a=1$, $b=0.3$ y $x_0=(2,0.8)$ ($V=1.736$), la solución cruza el eje $x_2=0$ doce veces en $[0,40]$, y $b\int_0^{40}x_2^2\dif t=V(x_0)-V(x(40))$ hasta la precisión del integrador (\cref{fig:iv-pendulo}).
\emph{¿Por qué no es global?} En $\R^2$ hay otros equilibrios $(k\pi,0)$, que están en $\{\Vdot=0\}$ y son soluciones que se quedan allí; además, $\{V\le c\}$ no es acotado en $\R^2$ y $V$ no es definida positiva fuera de $\abs{x_1}<2\pi$.
\end{ejemplo}

\begin{figure}[ht]
\centering
\includegraphics[width=\linewidth]{cap08/f_pendulo}
\caption{Péndulo con fricción ($a=1$, $b=0.3$). Izquierda: en azul, $\Nivel c$ ($c=1.8$), la componente de $\{V\le c\}$ que contiene al origen; en gris, las otras componentes; en coral, $E$; en violeta, el campo sobre $E$, vertical y transversal salvo en el origen; en negro, el nivel $V=2a$ de las separatrices. Los cuadrados son los cruces con $E$. Derecha: $V(x(t))$ y $\Vdot(x(t))$; en los cruces (líneas punteadas) $\Vdot=0$ y $V$ sigue bajando.}
\label{fig:iv-pendulo}
\end{figure}

\codigorepo{cap08_invarianza/C8_Ej2_PenduloLaSalle.m}{pendulo}{$\Nivel c$, $E$, los cruces y el balance $b\int x_2^2=V(x_0)-V(x(t))$}
\codigorepo{cap08_invarianza/C8_Ej2_PenduloLaSalle.m}{permanencia}{La prueba de permanencia con derivadas de Lie numéricas}

\video{V8_1_pendulo_lasalle}{Pasar por $E$ no es quedarse en $E$}

\begin{profundiza}[El péndulo vive en un cilindro]
El ángulo $x_1$ es un punto de la circunferencia $\mathbb S^1$, y el espacio de estado es el cilindro $\mathbb S^1\times\R$. Allí $V$ está bien definida y sus conjuntos de nivel son \emph{compactos} para todo $c$. El \cref{thm:iv-lasalle} se aplica globalmente con $M=\{(0,0),(\pi,0)\}$: \emph{toda} solución converge al equilibrio inferior o al superior. El superior es una silla cuya variedad estable es una curva, así que casi todas las condiciones iniciales terminan abajo. Pero ningún resultado puede dar estabilidad asintótica global: la región de atracción de un equilibrio asintóticamente estable es contráctil~\citep{BhSz70}, y el cilindro no lo es. Ningún campo localmente Lipschitz en el cilindro tiene un equilibrio globalmente asintóticamente estable, por razones topológicas.
\end{profundiza}

\begin{pausa}
Para $\dot x_1=x_2$, $\dot x_2=-x_1^3-x_2^3$, ¿qué concluye LaSalle con $V=\tfrac14x_1^4+\tfrac12x_2^2$? ¿Qué dice la linealización?
\respuesta{$\Vdot=x_1^3x_2+x_2(-x_1^3-x_2^3)=-x_2^4\le0$, $S=\{x_2=0\}$. Si $x_2\equiv0$, $\dot x_2=-x_1^3\equiv0$, luego $x_1\equiv0$: $M=\{0\}$. $V$ es radialmente no acotada: por el \cref{cor:iv-gas}, el origen es globalmente asintóticamente estable. La jacobiana en el origen es la matriz nula: el método indirecto no dice nada (\cref{tr:ly-imaginario}).}
\end{pausa}

%% ---------------------------------------------------------------------
\section{Qué puede ser $M$}
\label{sec:iv-ejemplos}
%% ---------------------------------------------------------------------

En el péndulo $M$ fue un punto. Los ejemplos de esta sección muestran las otras formas que toma $M$ y, con cada una, una trampa distinta.

\subsection{Un punto, pero en la componente correcta}

\begin{ejemplo}[label=ej:iv-duffing]{Resorte con amortiguamiento cuadrático \citep[Ej.~2.35]{KeBr23}}
Para $\dot x_1=x_2$, $\dot x_2=-(k_0x_1+k_1x_1^3)-\beta\,x_2\abs{x_2}$ con $k_0,\beta>0$, estudie el origen con la energía $V=\tfrac{k_0}2x_1^2+\tfrac{k_1}4x_1^4+\tfrac12x_2^2$ para un resorte duro ($k_1>0$) y uno blando ($k_1<0$).
\solucion $\Vdot=(k_0x_1+k_1x_1^3)x_2+x_2\bigl(-(k_0x_1+k_1x_1^3)-\beta x_2\abs{x_2}\bigr)=-\beta\abs{x_2}^3\le0$, con $S=\{x_2=0\}$. Si $x_2\equiv0$, $\dot x_2\equiv0$ obliga $x_1(k_0+k_1x_1^2)\equiv0$.
\emph{Resorte duro.} La única raíz real es $x_1=0$: $M=\{0\}$. $V$ es definida positiva y radialmente no acotada, y por el \cref{cor:iv-gas} el origen es \clave{globalmente asintóticamente estable}. El método indirecto no decide: la linealización es un centro, porque $\beta x_2\abs{x_2}$ es de segundo orden. La convergencia es lenta: la amplitud decae aproximadamente como $1/t$ (con $k_0=k_1=1$, $\beta=0.4$ y $x_0=(2,0)$, el producto amplitud $\times\,t$ vale $\approx6$ entre $t=100$ y $t=400$). El origen es globalmente asintóticamente estable y \emph{no} exponencialmente estable.
\emph{Resorte blando.} Aparecen dos sillas en $(\pm\sqrt{k_0/\abs{k_1}},0)$, y $V$ deja de ser radialmente no acotada (ni siquiera es definida positiva en $\R^2$). Sobre el eje $x_2=0$, $V$ alcanza su máximo en las sillas, con $c^*=k_0^2/(4\abs{k_1})$. Para $c<c^*$, la componente $\Nivel c$ que contiene al origen está en $\abs{x_1}<\sqrt{k_0/\abs{k_1}}$, es compacta y positivamente invariante, y $M=\{0\}$: el origen es \clave{localmente} asintóticamente estable y $\Nivel c$ está en su región de atracción.
\end{ejemplo}

\begin{trampa}[label=tr:iv-componente]{``$V(x_0)<c$, luego $x_0$ está en el estimado $\Nivel c$''}
\plausible{$\Nivel c$ es el conjunto donde $V\le c$, y $x_0$ cumple esa desigualdad.}
Con $k_0=1$, $k_1=-1$ y $\beta=0.4$, el punto $x_0=(1.2,0)$ tiene $V(x_0)=0.2016<c^*=0.25$ y \alerta{escapa a infinito} ($\abs{x_1}=50$ en $t=2.40$): está en otra componente de $\{V\le c\}$, la que rodea a las sillas por fuera (zona gris de la \cref{fig:iv-duffing}). El punto $(0.9,0)$, con $V=0.241$, sí está en $\Nivel c$ y converge.
\correcto{El estimado es la \emph{componente conexa} de $\{V\le c\}$ que contiene al origen (\cref{lem:iv-subnivel}). Es la \cref{tr:iv-subnivel} con consecuencias.}
\end{trampa}

\begin{figure}[ht]
\centering
\includegraphics[width=\linewidth]{cap08/f_duffing}
\caption{Izquierda: resorte duro ($k_1=1$), globalmente asintóticamente estable; niveles de $V$ en azul. Derecha: resorte blando ($k_1=-1$). En azul, $\Nivel c$ con $c=0.245$; en gris, la parte de $\{V\le c\}$ que \emph{no} es la componente del origen; en negro, el nivel $V=c^*=\tfrac14$ de las sillas. La solución coral arranca con $V(x_0)<c$ y escapa; la verde arranca en $\Nivel c$ y converge.}
\label{fig:iv-duffing}
\end{figure}

\codigorepo{cap08_invarianza/C8_Ej3_Duffing.m}{duffing}{Resorte duro y blando: $M=\{0\}$ y la componente equivocada}

\subsection{Un ciclo límite, con una $V$ que no es definida positiva}

\begin{ejemplo}[label=ej:iv-ciclo]{LaSalle prueba que un ciclo atrae}
Para $\dot x_1=x_1(1-x_1^2-x_2^2)-x_2$, $\dot x_2=x_2(1-x_1^2-x_2^2)+x_1$, que en polares es $\dot r=r(1-r^2)$, $\dot\theta=1$, pruebe que el círculo unitario atrae a las soluciones de un anillo.
\solucion Tome $V=\tfrac14(r^2-1)^2$. Como $\frac{\dif}{\dif t}r^2=2r^2(1-r^2)$,
\[
  \Vdot=\tfrac12(r^2-1)\cdot2r^2(1-r^2)=-r^2(1-r^2)^2\le0 .
\]
$V$ \emph{no} es definida positiva respecto a ningún punto: $V(0)=\tfrac14$ es un máximo local y $V=0$ en todo el círculo (el ``sombrero'' de la \cref{fig:iv-ciclo}). Para $0<c<\tfrac14$,
\[
  \Nivel c=\{V\le c\}=\bigl\{1-2\sqrt c\le r^2\le1+2\sqrt c\bigr\}
\]
es un anillo compacto que \emph{no contiene al origen} (con $c=0.2$, $0.325\le r\le1.376$) y es positivamente invariante porque $\Vdot\le0$ en todo $\R^2$. $E=\Nivel c\cap\{r=1\}$ es el círculo, que es una órbita periódica ($r\equiv1\implica\dot r=0$, $\dot\theta=1$) y por lo tanto invariante: $M=E$. Por el \cref{thm:iv-lasalle}, \clave{toda solución que arranca en el anillo tiende al ciclo límite}.
\end{ejemplo}

Tres comentarios. Poincaré--Bendixson (\cref{sec:oc-pb}) da la \emph{existencia} de una órbita periódica en el anillo; LaSalle da su \emph{atracción}. Si se hubieran revisado solo los equilibrios de $E$ (no hay ninguno), se habría concluido $M=\emptyset$ (\cref{tr:iv-equilibrios}). Y si $c\ge\tfrac14$, $\Nivel c$ contiene al origen, $E=\{0\}\cup\{r=1\}$, y LaSalle solo da convergencia a esa unión: para descartar el origen hace falta un argumento adicional (aquí, que $r(t)$ es monótona).

\begin{figure}[ht]
\centering
\includegraphics[width=\linewidth]{cap08/f_ciclo}
\caption{Izquierda: el anillo $\Nivel{0.2}$ (azul), el conjunto $E=M$ (coral) y dos soluciones que convergen a él. Derecha: el perfil de $V=\tfrac14(r^2-1)^2$ y las mismas soluciones, muestreadas en el tiempo, ``rodando'' hacia el fondo del sombrero.}
\label{fig:iv-ciclo}
\end{figure}

\codigorepo{cap08_invarianza/C8_Ej4_CicloLimite.m}{ciclo}{El anillo $\Nivel c$ y la convergencia al ciclo}

\video{V8_2_ciclo_sombrero}{Rodar por el sombrero hasta el ciclo}

\subsection{Una recta de equilibrios}

\begin{ejemplo}[label=ej:iv-adaptativo]{Control adaptativo \citep[Ej.~4.10]{Kha02}, \citep{Ser00}}
Para la planta $\dot y=ay+u$ con $a$ \emph{desconocido}, la ley adaptativa $u=-ky$, $\dot k=\gamma y^2$ ($\gamma>0$) da el lazo
\[
  \dot y=-(k-a)\,y,\qquad \dot k=\gamma\,y^2 .
\]
Pruebe que $y(t)\to0$.
\solucion La recta $y=0$ está hecha de equilibrios: el origen no es un equilibrio aislado y no puede ser asintóticamente estable. Para cualquier constante $b>a$, tome $V=\tfrac12y^2+\frac1{2\gamma}(k-b)^2$:
\[
  \Vdot=-(k-a)\,y^2+(k-b)\,y^2=-(b-a)\,y^2\le0 .
\]
$V$ es radialmente no acotada en $(y,k)$, así que $\Nivel c=\{V\le c\}$ es compacto y positivamente invariante para todo $c$. $E=\Nivel c\cap\{y=0\}$, y como cada punto de $E$ es un equilibrio, $M=E$. Por el \cref{thm:iv-lasalle}, $y(t)\to0$ para toda condición inicial. La constante $b$ solo se usa en el análisis: basta que \emph{exista}, y el controlador no necesita conocer $a$.
\end{ejemplo}

LaSalle da convergencia a la recta. En este caso hay convergencia a un punto, por un argumento adicional: $k(t)$ es no decreciente y acotado (está en $\Nivel c$), luego converge a algún $k_\infty$. Con $b=a$, $W=\tfrac12y^2+\frac1{2\gamma}(k-a)^2$ cumple $\dot W=0$: es una cantidad \emph{conservada}, y con $\gamma=1$ las trayectorias son arcos de circunferencias centradas en $(0,a)$. Como $k$ crece y $y\to0$, si $y_0\neq0$ el arco termina en la intersección superior con el eje:
\begin{equation}
  k_\infty=a+\sqrt{(k_0-a)^2+\gamma\,y_0^2}\;>\;a .
  \label{eq:iv-kinf}
\end{equation}
Con $a=\gamma=1$ y $(y_0,k_0)=(2,-1)$, $k_\infty=1+2\sqrt2=3.83$ (\cref{fig:iv-adaptativo}).

\begin{trampa}[label=tr:iv-punto]{``$x(t)\to M$, luego $x(t)$ converge a un punto de $M$, y ese punto es el que usé en $V$''}
\plausible{LaSalle se aplicó con $V$ centrada en $(0,b)$, así que la solución termina en $(0,b)$.}
El límite \eqref{eq:iv-kinf} depende de la condición inicial y \emph{no es $b$}: $b$ es una ficción del análisis. Que la solución converja a un \emph{punto} de $M$ requirió un argumento aparte (la monotonía de $k$); en el \cref{ej:iv-ciclo}, la solución no converge a ningún punto de $M$. Observe además que con $b=a$ se tiene $\Vdot\equiv0$ y LaSalle no da información: escoger $b>a$ es lo que hace funcionar el argumento.
\correcto{LaSalle da convergencia al \emph{conjunto} $M$. La convergencia a un punto, y cuál, requieren argumentos adicionales.}
\end{trampa}

\begin{figure}[ht]
\centering
\includegraphics[width=\linewidth]{cap08/f_adaptativo}
\caption{Izquierda: plano $(y,k)$ con $a=\gamma=1$. Los equilibrios con $k<a$ (coral) son inestables y los con $k>a$ (verde) estables. Las curvas punteadas son niveles de $V$ con $b=2$; las estrellas, los $k_\infty$ de \eqref{eq:iv-kinf}. Derecha: $y(t)$ y $k(t)$ para $(y_0,k_0)=(2,-1)$; $k_\infty=3.83\neq b$.}
\label{fig:iv-adaptativo}
\end{figure}

\codigorepo{cap08_invarianza/C8_Ej5_Adaptativo.m}{adaptativo}{$k(t)$ simulado contra la fórmula de $k_\infty$}

\subsection{Un subespacio no observable}

\begin{ejemplo}[label=ej:iv-observable]{La prueba del \cref{thm:ly-observable}}
Sea $\dot x=Ax$ con $P=P\T>0$ tal que $A\T P+PA=-C\T C$. Pruebe que si $(A,C)$ es observable, $A$ es Hurwitz, y describa $M$ cuando no lo es.
\solucion $V=x\T Px$ es definida positiva y radialmente no acotada, con $\Vdot=-\norm{Cx}^2\le0$ y $S=\{x:\ Cx=0\}$. La prueba de permanencia con $h(x)=Cx$: si una solución se queda en $S$, entonces $y(t)=Ce^{At}x_0\equiv0$, y derivando en $t=0$,
\[
  Cx_0=0,\ CAx_0=0,\ \dots,\ CA^{n-1}x_0=0\iff\mathcal O\,x_0=0,\qquad \mathcal O=\begin{bmatrix}C\\CA\\\vdots\\CA^{n-1}\end{bmatrix}.
\]
Si $(A,C)$ es observable, $\mathcal O$ tiene rango completo y $x_0=0$: $M=\{0\}$, y por el \cref{cor:iv-gas} el origen es globalmente asintóticamente estable, es decir, $A$ es Hurwitz. \alerta{La prueba de permanencia es literalmente la prueba de observabilidad, y $M$ es el subespacio no observable} $\ker\mathcal O$.
Las dos masas del \cref{cap:lyapunov} (paredes con resortes unitarios, acople $\kappa$, amortiguador $b=0.3$ solo en la masa~1) lo muestran. Con $x=(q_1,v_1,q_2,v_2)$ y la energía $V=\tfrac12\bigl(q_1^2+q_2^2+\kappa(q_1-q_2)^2+v_1^2+v_2^2\bigr)$, $\Vdot=-0.3\,v_1^2$. Con $\kappa=0.4$: $v_1\equiv0$ obliga $\dot v_1=-q_1-\kappa(q_1-q_2)\equiv0$, y como $\dot q_1=v_1=0$, $q_1$ y $q_2$ son constantes; entonces $v_2=\dot q_2=0$ y $\dot v_2=-q_2+\kappa(q_1-q_2)=0$, lo que junto con la primera relación da $q_1=q_2=0$. El rango de observabilidad es $4$ y $M=\{0\}$. Con $\kappa=0$, el rango es $2$ y $M=\ker\mathcal O=\{q_1=v_1=0\}$: la masa~2 oscila para siempre en un plano invariante que la disipación no ve (\cref{fig:ly-masas}).
\end{ejemplo}

\codigorepo{cap08_invarianza/C8_Ej6_Observabilidad.m}{observabilidad}{$M=\ker\mathcal O$ para las dos masas}

\subsection{Un punto interior: Lotka--Volterra}

\begin{ejemplo}[label=ej:iv-lotka]{La función de Volterra}
Con la notación de la \cref{sec:oc-lotka} ($v$: presas, $p$: depredadores) y autolimitación solo de las presas,
\[
  \dot v=v\,(a-ev-bp),\qquad \dot p=p\,(-c+dv),\qquad e>0,\ ad>ec ,
\]
pruebe que el equilibrio de coexistencia $v^*=c/d$, $p^*=(a-ev^*)/b>0$ atrae todo el cuadrante abierto.
\solucion En $\D=\{v>0,\ p>0\}$ tome la función de Volterra
\[
  V(v,p)=d\Bigl(v-v^*-v^*\ln\frac v{v^*}\Bigr)+b\Bigl(p-p^*-p^*\ln\frac p{p^*}\Bigr).
\]
Como $z-1-\ln z\ge0$ con igualdad solo en $z=1$, $V\ge0$ y $V=0$ solo en $(v^*,p^*)$. Usando $a=ev^*+bp^*$ y $c=dv^*$,
\[
  \Vdot=d(v-v^*)\bigl(-e(v-v^*)-b(p-p^*)\bigr)+b(p-p^*)\,d(v-v^*)=-de\,(v-v^*)^2\le0 .
\]
$V\to\infty$ cuando $v\to0^+$, $p\to0^+$ o $\norm{(v,p)}\to\infty$, así que los conjuntos $\{V\le\ell\}$, $\ell>0$, son compactos \emph{dentro de} $\D$ y positivamente invariantes. $E=\{v=v^*\}$, y si $v\equiv v^*$, $\dot v\equiv0$ da $a-ev^*-bp\equiv0$, es decir $p\equiv p^*$: $M=\{(v^*,p^*)\}$. El equilibrio es asintóticamente estable y su región de atracción es \emph{todo el cuadrante abierto}. Con $a=b=d=1$, $c=0.5$ y $e=0.2$: $(v^*,p^*)=(0.5,\,0.9)$. Con $e=0$, $\Vdot\equiv0$: $V$ es la cantidad conservada del modelo clásico y las órbitas son cerradas (\cref{fig:iv-lotka}).
\end{ejemplo}

Este es el argumento que la \cref{sec:oc-lotka} dejó pendiente. No es estabilidad asintótica global en el cuadrante \emph{cerrado}: los ejes son invariantes y las soluciones que nacen en ellos no llegan a $(v^*,p^*)$.

\begin{figure}[ht]
\centering
\includegraphics[width=\linewidth]{cap08/f_lotka}
\caption{Lotka--Volterra con $a=b=d=1$, $c=0.5$. Izquierda: $e=0$, la función de Volterra se conserva y la órbita es cerrada. Derecha: $e=0.2$, la solución cruza transversalmente $E=\{v=v^*\}$ y converge a $(v^*,p^*)$.}
\label{fig:iv-lotka}
\end{figure}

\codigorepo{cap08_invarianza/C8_Ej7_LotkaVolterra.m}{lotka}{La función de Volterra, conservada y disipada}

\subsection{Convergencia sin estabilidad}

\begin{ejemplo}[label=ej:iv-inestable]{LaSalle da convergencia, no estabilidad}
En polares, sea $\dot r=r(1-r)$, $\dot\theta=1-\cos\theta\ (\ge0)$, en el anillo $\Omega=\{\tfrac12\le r\le\tfrac32\}$, donde el campo es suave. ¿Qué dice LaSalle con $V=(1-r)^2$? ¿Es estable el equilibrio $(1,0)$?
\solucion $\Vdot=-2(1-r)\,r(1-r)=-2r(1-r)^2\le0$, y $\Omega=\{V\le\tfrac14\}$ (que excluye al origen, donde $V=1$) es compacto y positivamente invariante. $E=\{r=1\}$, el círculo es invariante y $M=E$: LaSalle concluye que $x(t)$ tiende al círculo. La dinámica real es más fina. Sobre el círculo, $\theta$ crece hasta $2\pi\equiv0$, y toda solución del anillo converge al equilibrio $(1,0)$. Pero $(1,0)$ es \clave{inestable}: una condición inicial a distancia $0.15$ del equilibrio, en $\theta=0.15$, da la vuelta completa y se aleja hasta distancia $2$ antes de regresar (\cref{fig:iv-inestable}).
\end{ejemplo}

\begin{trampa}[label=tr:iv-estabilidad]{``Toda solución tiende a $\eq$, luego $\eq$ es asintóticamente estable''}
\plausible{LaSalle mostró que todas las soluciones convergen, y en eso consiste la estabilidad asintótica.}
La estabilidad asintótica pide \emph{estabilidad} y atracción (\cref{tr:ly-convergencia}). En el \cref{ej:iv-inestable} hay atracción sin estabilidad. Los \cref{cor:iv-as,cor:iv-gas} exigen $V$ definida positiva \emph{precisamente} para obtener la estabilidad con el \cref{thm:ly-directo}.
\correcto{LaSalle solo da convergencia. La estabilidad hay que probarla aparte. Además, $M$ puede ser mucho más grande que los conjuntos límite reales: aquí $M$ es el círculo y todos los $\omegalim{x_0}$ son $\{(1,0)\}$.}
\end{trampa}

\begin{figure}[ht]
\centering
\includegraphics[width=\linewidth]{cap08/f_inestable}
\caption{Izquierda: el anillo $\Omega$, el conjunto $E=M$ y tres soluciones que convergen a $(1,0)$. Derecha: distancia a $(1,0)$; la solución negra arranca a distancia $0.15$ y hace una excursión de tamaño $2$.}
\label{fig:iv-inestable}
\end{figure}

\codigorepo{cap08_invarianza/C8_Ej8_AtractivoInestable.m}{inestable}{Atractivo pero inestable}

\begin{pausa}
En el \cref{ej:iv-adaptativo}, ¿qué pasa si $y_0=0$? ¿Y si $b<a$?
\respuesta{Si $y_0=0$, la solución es el equilibrio $(0,k_0)$ y $k_\infty=k_0$; la fórmula \eqref{eq:iv-kinf} da $a+\abs{k_0-a}$, que solo coincide si $k_0\ge a$: la fórmula supone $y_0\neq0$. Si $b<a$, $\Vdot=(a-b)y^2\ge0$ y LaSalle no se aplica con esa $V$; por eso se escoge $b>a$, que siempre existe.}
\end{pausa}

%% ---------------------------------------------------------------------
\section{Estimar la región de atracción cuando $\Vdot\le0$}
\label{sec:iv-region}
%% ---------------------------------------------------------------------

La \cref{prop:ly-estimacion} exigía $\Vdot<0$ fuera del origen. Con LaSalle basta $\Vdot\le0$, siempre que en el conjunto de nivel ninguna solución no trivial pueda quedarse donde $\Vdot=0$. Las dos versiones se escriben juntas así.

\begin{proposicion}[label=prop:iv-estimacion]{Estimado de la región de atracción con un conjunto de nivel}
Sea $V$ de clase $\Cont1$, definida positiva en un dominio $\D\ni0$, y $\Nivel c$ la componente conexa de $\{x\in\D:\ V(x)\le c\}$ que contiene al origen, acotada y con $\overline{\Nivel c}\subset\D$. Si
\begin{enumerate}
  \item[(a)] $\Vdot<0$ en $\Nivel c\setminus\{0\}$, o
  \item[(b)] $\Vdot\le0$ en $\Nivel c$ y ninguna solución distinta de $x\equiv0$ permanece en $\{x\in\Nivel c:\ \Vdot(x)=0\}$,
\end{enumerate}
entonces $\Nivel c\subset\RA$. En el caso (a), el mayor $c$ admisible es
\begin{equation}
  c^*=\inf\bigl\{V(x):\ x\in\D,\ x\neq0,\ \Vdot(x)\ge0\bigr\},
  \label{eq:iv-cestrella}
\end{equation}
y cualquier $c<c^*$ sirve.
\end{proposicion}

\begin{demostracion}
Por el \cref{lem:iv-subnivel}, $\Nivel c$ es compacto y positivamente invariante. En el caso (b), el \cref{thm:iv-lasalle} da $x(t)\to M=\{0\}$. En el caso (a), $E=\{0\}$. En los dos, $x(t)\to0$ desde todo punto de $\Nivel c$.
\end{demostracion}

\notamargen{En la práctica, $c^*$ se calcula analíticamente cuando se puede y, si no, en una malla: se evalúa $V$ en los puntos donde $\Vdot\ge0$ y se toma el mínimo. Es lo que hace \cmd{nivelCritico} (en \cmd{comun/}). La malla debe cubrir solo $\D$: en el péndulo inclinado del \cref{ej:ly-cpendulo}, una malla más ancha que $\abs{x_1}<\pi$ encuentra puntos con $\Vdot\ge0$ más allá de $\pm\pi$ y da un $c^*$ menor que $2.308$, por la trampa de las componentes.}

\begin{ejemplo}[label=ej:iv-region]{El problema que el capítulo~7 dejó a medias}
Para $\dot x_1=x_2$, $\dot x_2=-x_1+x_1^3-x_2$ (problema~\ref{prob:ly-duffing}), estime la región de atracción del origen con $V=\tfrac12x_1^2-\tfrac14x_1^4+\tfrac12x_2^2$.
\solucion $\Vdot=(x_1-x_1^3)x_2+x_2(-x_1+x_1^3-x_2)=-x_2^2\le0$: solo semidefinida, y el caso (a) no se aplica. Las sillas $(\pm1,0)$ tienen $V=\tfrac14$, y para $c<\tfrac14$ la componente $\Nivel c$ que contiene al origen está en $\abs{x_1}<1$. \emph{Permanencia}: $x_2\equiv0$ obliga $x_1-x_1^3\equiv0$, es decir $x_1\in\{0,\pm1\}$, y en $\abs{x_1}<1$ solo queda $x_1=0$. Por el caso (b),
\[
  \bigl\{V<\tfrac14\bigr\}\cap\bigl\{\abs{x_1}<1\bigr\}\subset\RA .
\]
El punto $(0,0.7)$, con $V=0.245$, converge; los de la componente exterior (gris en la \cref{fig:iv-region}) pueden escapar aunque tengan $V<\tfrac14$.
\end{ejemplo}

\begin{figure}[ht]
\centering
\includegraphics[width=0.66\linewidth]{cap08/f_region}
\caption{El estimado del \cref{ej:iv-region} (azul), limitado por las separatrices de las sillas en el nivel $V=\tfrac14$. En gris, la parte de $\{V<\tfrac14\}$ fuera de $\abs{x_1}<1$. Soluciones azules: convergen; corales: escapan.}
\label{fig:iv-region}
\end{figure}

\codigorepo{cap08_invarianza/C8_Ej9_RegionLaSalle.m}{region}{El estimado con $\Vdot\le0$ y la función \cmd{nivelCritico}}

\begin{trampa}[label=tr:iv-estimado]{``Fuera de $\Nivel{c^*}$ las soluciones no convergen''}
\plausible{$c^*$ es el mayor nivel que se puede certificar; si se pudiera más, la proposición lo daría.}
$c^*$ es el mayor nivel que se puede certificar \emph{con esa $V$}. Los estimados son conservadores: en el \cref{ej:iv-g1} todas las soluciones de prueba convergen, también desde fuera de $\Nivel{c^*}$; en el \cref{ej:iv-z3}, la mejor elipse cubre el 47\,\% de la región real.
\correcto{$\Nivel{c^*}\subset\RA$ es una inclusión, no una igualdad. Para mejorar el estimado hay que cambiar de $V$: es el tema del resto del capítulo.}
\end{trampa}

%% ---------------------------------------------------------------------
\section{Construir $V$: el método del gradiente variable}
\label{sec:iv-gradiente}
%% ---------------------------------------------------------------------

Los teoremas de Lyapunov y el principio de invarianza dicen qué hacer \emph{una vez que se tiene} una función $V$. Los teoremas conversos (\cref{sec:av-conversos}) garantizan que, si el origen es asintóticamente estable, una tal $V$ existe, pero no dicen cómo encontrarla. En sistemas físicos la energía es la primera candidata. Cuando no hay energía evidente, o cuando la energía solo da $\Vdot\le0$, hacen falta métodos sistemáticos. Se verán tres, y cada uno \emph{fija} algo distinto y deja libre el resto:

\begin{center}
\small
\begin{tabular}{@{}p{2.6cm}p{3.6cm}p{5.4cm}@{}}
\toprule
Método & Qué se fija & Qué se obtiene \\
\midrule
Gradiente variable & la forma de $\nabla V$ & $V$ y un estimado de $\RA$, con parámetros libres para mejorarlo \\
Krasovskii & $V=f\T Pf$ & una condición sobre la jacobiana; estabilidad global si vale en $\R^n$ \\
Zubov & la forma de $\Vdot$ (una EDP) & la región de atracción \emph{exacta} \\
\bottomrule
\end{tabular}
\end{center}

La idea del gradiente variable es no adivinar $V$, sino su \clave{gradiente} $g(x)=\nabla V(x)$, porque $\Vdot$ depende de $V$ solo a través de él: $\Vdot(x)=g(x)\T f(x)$. Se propone $g$ con coeficientes libres, se escogen para que $g\T f$ sea definida negativa (o semidefinida) y al final se recupera $V$ integrando. El punto delicado es que no cualquier campo $g$ es un gradiente.

\begin{lema}[label=lem:iv-gradiente]{Condición de gradiente}
Sea $g:\D\to\R^n$ de clase $\Cont1$ en un dominio $\D$ convexo (o estrellado respecto al origen). Existe $V$ con $\nabla V=g$ si y solo si la jacobiana de $g$ es simétrica:
\[
  \jac{g_i}{x_j}=\jac{g_j}{x_i},\qquad i,j=1,\dots,n .
\]
En ese caso $V(x)=\int_0^1g(\sigma x)\T x\dif\sigma$, la integral de línea $\int_0^xg\T\dif y$ no depende del camino, y se puede integrar a lo largo de los ejes:
\[
  V(x)=\int_0^{x_1}g_1(y_1,0,\dots,0)\dif y_1+\int_0^{x_2}g_2(x_1,y_2,0,\dots,0)\dif y_2+\cdots+\int_0^{x_n}g_n(x_1,\dots,x_{n-1},y_n)\dif y_n .
\]
\end{lema}

La necesidad es la igualdad de las derivadas cruzadas de $V$ (teorema de Schwarz); la suficiencia en un dominio estrellado es el lema de Poincaré.

\begin{trampa}[label=tr:iv-simetria]{``Integré $g$ por los ejes, luego $\nabla V=g$''}
\plausible{La fórmula del lema da una función; su gradiente debe ser el $g$ con el que se construyó.}
Si $g$ no es simétrico, la integral por los ejes produce \emph{alguna} función, pero su gradiente no es $g$, y la cuenta $\Vdot=g\T f$ es falsa. Con $g=(x_1,\ x_1+x_2)$, $\partial g_1/\partial x_2=0\neq1=\partial g_2/\partial x_1$. Integrando por los ejes, $V=\tfrac12x_1^2+x_1x_2+\tfrac12x_2^2=\tfrac12(x_1+x_2)^2$, cuyo gradiente es $(x_1+x_2,\ x_1+x_2)\neq g$; además $V$ ni siquiera es definida positiva.
\correcto{Primero se impone la simetría y después se integra.}
\end{trampa}

\begin{receta}[El gradiente variable (Schultz y Gibson)]
\begin{enumerate}
  \item Proponer $g(x)$ con componentes $g_i=a_{i1}x_1+\cdots+a_{in}x_n$, con $a_{ij}$ constantes o funciones de $x$.
  \item Calcular $\Vdot=g\T f$.
  \item Escoger los $a_{ij}$ para que $\Vdot$ sea definida negativa (típicamente, anulando términos cruzados sin signo).
  \item Imponer la simetría $\partial g_i/\partial x_j=\partial g_j/\partial x_i$.
  \item Integrar para obtener $V$ y verificar que sea definida positiva; si no lo es, volver al paso~3.
  \item Estimar la región de atracción con la \cref{prop:iv-estimacion}.
\end{enumerate}
\end{receta}

\begin{ejemplo}[label=ej:iv-g1]{El ejercicio del handout, resuelto}
Estudie el origen de
\begin{equation}
  \dot x_1=x_2,\qquad \dot x_2=-(x_1+x_2)-\sin(x_1+x_2) .
  \label{eq:iv-g1}
\end{equation}
\solucion El único equilibrio es el origen: $x_2=0$ y $x_1+\sin x_1=0$ obligan $x_1=0$.
\emph{Pasos 1 y 4.} Con $g=(a_{11}x_1+a_{12}x_2,\ a_{21}x_1+a_{22}x_2)$ y coeficientes constantes, la simetría pide $a_{12}=a_{21}$. Multiplicar $g$ por una constante positiva solo escala $V$; se toma $a_{12}=a_{21}=a_{22}=1$ y queda un parámetro libre: $g=(a_{11}x_1+x_2,\ x_1+x_2)$.
\emph{Paso 2.} Con $y=x_1+x_2$,
\[
  \Vdot=(a_{11}x_1+x_2)\,x_2+y\,(-y-\sin y)=(a_{11}-2)\,x_1x_2-x_1^2-y\sin y .
\]
\emph{Paso 3.} El término $x_1x_2$ no tiene signo; se anula con $a_{11}=2$, y $\Vdot=-x_1^2-y\sin y$. Como $y\sin y>0$ si $0<\abs y<\pi$, $\Vdot$ es \clave{definida negativa en la franja} $\{\abs{x_1+x_2}<\pi\}$.
\emph{Paso 5.} Integrando por los ejes,
\[
  V=\int_0^{x_1}2y_1\dif y_1+\int_0^{x_2}(x_1+y_2)\dif y_2=x_1^2+x_1x_2+\tfrac12x_2^2=\tfrac12\bigl(x_1^2+(x_1+x_2)^2\bigr),
\]
definida positiva y radialmente no acotada. El origen es asintóticamente estable.
\emph{Paso 6.} En las coordenadas $(x_1,y)$, $V=\tfrac12(x_1^2+y^2)$ tiene niveles circulares, y $\Vdot\ge0$ solo es posible con $\abs y\ge\pi$. El mínimo de $V$ sobre $\abs y=\pi$ es $\pi^2/2$, en $x_1=0$: $c^*=\pi^2/2$ y
\[
  \bigl\{x:\ x_1^2+(x_1+x_2)^2<\pi^2\bigr\}\subset\RA ,
\]
una elipse de área $\pi^3=31.0$ (la transformación $(x_1,x_2)\mapsto(x_1,y)$ tiene determinante $1$).
\emph{¿Es global?} Con esta $V$ no se puede concluir: en las bandas $\pi<\abs y<2\pi$ con $x_1$ pequeño, $\Vdot>0$ (zonas coral de la \cref{fig:iv-gradiente}). Sin embargo, las $49$ soluciones de prueba desde una malla de $[-12,12]^2$ convergen, incluso las que arrancan lejos del estimado.
\end{ejemplo}

\begin{figure}[ht]
\centering
\includegraphics[width=\linewidth]{cap08/f_gradiente}
\caption{\Cref{ej:iv-g1}. Izquierda ($a_{11}=2$): en coral, la región donde $\Vdot>0$; punteadas, $x_1+x_2=\pm\pi$; en azul, la frontera del estimado; soluciones negras (dentro) y grises (fuera), todas convergen. Derecha: área del estimado $\Nivel{c^*}$ en función de $a_{11}$.}
\label{fig:iv-gradiente}
\end{figure}

\paragraph{¿Por qué $a_{11}=2$?} Con $a_{11}$ libre, $V=\tfrac{a_{11}}2x_1^2+x_1x_2+\tfrac12x_2^2$ es definida positiva si y solo si $a_{11}>1$. Cerca del origen, con $y\sin y\approx y^2$, la parte cuadrática de $\Vdot$ es
\[
  -2x_1^2+(a_{11}-4)\,x_1x_2-x_2^2 ,
\]
definida negativa si y solo si $(a_{11}-4)^2<8$, es decir $a_{11}\in(4-2\sqrt2,\ 4+2\sqrt2)\approx(1.17,\ 6.83)$. \emph{Cualquier} $a_{11}$ en ese intervalo prueba la estabilidad asintótica local. La elección $a_{11}=2$ es la que hace el álgebra más limpia, no la mejor: el mayor estimado está en $a_{11}\approx1.48$, con área $37.4$ contra $31.0$ (\cref{fig:iv-gradiente,fig:iv-barrido}).

\begin{figure}[ht]
\centering
\includegraphics[width=\linewidth]{cap08/f_barrido}
\caption{Barrido de $a_{11}$ en el \cref{ej:iv-g1}. En coral, la región donde $\Vdot>0$; en azul, la frontera del mayor $\{V\le c\}$ que no la toca. Con $a_{11}=1.1$ el estimado colapsa: $\Vdot$ no es definida negativa cerca del origen.}
\label{fig:iv-barrido}
\end{figure}

\codigorepo{cap08_invarianza/C8_Ej10_GradienteVariable.m}{g1}{$V$, $\Vdot$ y el estimado con $a_{11}=2$}
\codigorepo{cap08_invarianza/C8_Ej10_GradienteVariable.m}{barrido}{El área del estimado como función de $a_{11}$}

\video{V8_3_gradiente_barrido}{Mover $a_{11}$ y ver crecer (y colapsar) el estimado}

\begin{trampa}[label=tr:iv-limpia]{``La elección que anula el término cruzado es la mejor''}
\plausible{Con $a_{11}=2$ desaparece el término sin signo y $\Vdot$ queda lo más negativa posible.}
Escoger los parámetros para que $\Vdot$ ``se vea bonita'' es una heurística algebraica, no un criterio de optimalidad. Con $a_{11}=1.48$ queda un término cruzado en $\Vdot$ y el estimado es un 20\,\% más grande.
\correcto{Si el objetivo es el mayor estimado de la región de atracción, hay que optimizar sobre los parámetros libres, en la práctica numéricamente (\cmd{C8\_Ej10}).}
\end{trampa}

\subsection{De dónde salió la función inclinada del péndulo}

En el \cref{ej:ly-cruzado} se \emph{propuso} una $V$ con un término cruzado para el péndulo con fricción, y resultó estricta. El gradiente variable la \emph{deduce}.

\begin{ejemplo}[label=ej:iv-g2]{El gradiente variable encuentra la $V$ del capítulo~7}
Para $\dot x_1=x_2$, $\dot x_2=-\sin x_1-b\,x_2$ ($b>0$), busque una $V$ con $\Vdot$ definida negativa.
\solucion \emph{Pasos 1 a 4.} Proponga $g_2=k\,x_1+x_2$, con $k$ por determinar. La simetría pide $\partial g_1/\partial x_2=\partial g_2/\partial x_1=k$, así que $g_1=\varphi(x_1)+k\,x_2$ para alguna $\varphi$. Entonces
\[
  \Vdot=\bigl(\varphi(x_1)+k\,x_2\bigr)x_2+(k\,x_1+x_2)(-\sin x_1-b\,x_2)=\bigl(\varphi(x_1)-\sin x_1-kb\,x_1\bigr)x_2+(k-b)\,x_2^2-k\,x_1\sin x_1 .
\]
El primer término es lineal en $x_2$, sin signo; se anula con $\varphi(x_1)=\sin x_1+kb\,x_1$, y queda
\[
  \Vdot=-k\,x_1\sin x_1-(b-k)\,x_2^2 ,
\]
definida negativa en $\{\abs{x_1}<\pi\}$ si $0<k<b$.
\emph{Paso 5.} Integrando,
\[
  V=(1-\cos x_1)+\tfrac12\bigl(kb\,x_1^2+2k\,x_1x_2+x_2^2\bigr),
\]
que es \eqref{eq:ly-Vcruzado} con $a=1$ y $p_{12}=k$. La forma cuadrática es definida positiva si y solo si $kb-k^2>0$: la misma condición $k<b$. Con $k\to0$ se recupera la energía.
\emph{Paso 6.} El mínimo de $V$ sobre $\abs{x_1}=\pi$ es $c^*(k)=2+\frac{\pi^2}2k(b-k)$, máximo en $k=b/2$; con $b=0.5$, $c^*=2.308$, el valor del \cref{ej:ly-cpendulo}.
\end{ejemplo}

El término cruzado ``inclina'' los niveles de $V$ (\cref{fig:iv-gradiente-pendulo}). Por eso $\Vdot$ es negativa también sobre el eje $x_2=0$, donde la energía se quedaba quieta, y basta el método directo, sin LaSalle. Tener una $V$ \emph{estricta} importa: los argumentos de robustez del \cref{cap:avanzada} (conversos, UUB, ISS) necesitan $\Vdot\le-\alpha_3(\norm x)$, no solo $\Vdot\le0$.

\begin{figure}[ht]
\centering
\includegraphics[width=\linewidth]{cap08/f_gradiente_pendulo}
\caption{\Cref{ej:iv-g2} con $b=0.5$ y $k=b/2$. Izquierda: niveles de la energía (gris, punteados) y de la $V$ del gradiente variable (azul, inclinados); en verde, $\Nivel{c^*}$, que toca $\abs{x_1}=\pi$. Derecha: a lo largo de la misma solución, la energía baja con escalones y $V$ decrece estrictamente.}
\label{fig:iv-gradiente-pendulo}
\end{figure}

\codigorepo{cap08_invarianza/C8_Ej11_GradientePendulo.m}{pendulo}{Energía contra $V$ estricta y $c^*(k)$}

\begin{pausa}
¿Qué pasa en el \cref{ej:iv-g2} si se propone $g_2=k\,x_1+x_2$ con $k>b$?
\respuesta{$\Vdot=-kx_1\sin x_1+(k-b)x_2^2$ deja de ser semidefinida negativa (es positiva sobre el eje $x_1=0$), y además la forma cuadrática de $V$ deja de ser definida positiva ($kb-k^2<0$). Las dos condiciones fallan a la vez: no hay $V$ útil con esa forma.}
\end{pausa}

%% ---------------------------------------------------------------------
\section{El método de Krasovskii}
\label{sec:iv-krasovskii}
%% ---------------------------------------------------------------------

La idea de Krasovskii es medir la ``energía'' no del estado sino de la \clave{velocidad}. Con $z=f(x)$, a lo largo de las soluciones
\[
  \dot z=\jac fx(x)\,\dot x=J(x)\,z,\qquad J(x)=\jac fx(x) .
\]
La velocidad obedece un sistema ``lineal'' con una matriz $J(x(t))$ que cambia con el tiempo. Si una forma cuadrática $z\T Pz$ decrece para \emph{todas} las matrices $J(x)$ que pueden aparecer, la velocidad se apaga y el sistema se detiene; si además el único lugar donde $f=0$ es el origen, se detiene allí.\notamargen{Para $f(x)=Ax$, $J=A$ y la condición del teorema es $A\T P+PA\le-R$: la ecuación de Lyapunov. La candidata es $V=x\T A\T PA\,x$.}

\begin{teorema}[label=thm:iv-krasovskii]{Krasovskii \citep{Kra63}, \citep[ejercicios del cap.~4]{Kha02}}
Sea $x=0$ un equilibrio de $\dot x=f(x)$, con $f:\D\to\R^n$ de clase $\Cont1$ y $\D$ un dominio \clave{convexo} que contiene al origen. Si existen $P=P\T>0$ y $R=R\T>0$ con
\begin{equation}
  J(x)\T P+PJ(x)\le-R\qquad\text{para todo }x\in\D,
  \label{eq:iv-krasovskii}
\end{equation}
entonces: (i) el origen es el \clave{único equilibrio} en $\D$; (ii) el origen es asintóticamente estable, con función de Lyapunov $V(x)=f(x)\T Pf(x)$; (iii) si $\D=\R^n$, es globalmente asintóticamente estable.
\end{teorema}

\begin{demostracion}
(i) Como $\D$ es convexo, el segmento $\{\sigma x:\ \sigma\in[0,1]\}$ está en $\D$, y $f(x)=\int_0^1J(\sigma x)\,x\dif\sigma$. Entonces
\[
  x\T Pf(x)=\int_0^1x\T PJ(\sigma x)\,x\dif\sigma=\tfrac12\int_0^1x\T\bigl(PJ(\sigma x)+J(\sigma x)\T P\bigr)x\dif\sigma\le-\tfrac12\,x\T Rx .
\]
Por Cauchy--Schwarz, $\lmax(P)\norm x\norm{f(x)}\ge\tfrac12\lmin(R)\norm x^2$, es decir
\[
  \norm{f(x)}\ge\kappa\norm x,\qquad \kappa=\frac{\lmin(R)}{2\lmax(P)}>0 ;
\]
en particular $f(x)\neq0$ si $x\neq0$. (ii) $V=f\T Pf\ge\lmin(P)\kappa^2\norm x^2$ es definida positiva, y
\[
  \Vdot=\dot z\T Pz+z\T P\dot z=f\T\bigl(J\T P+PJ\bigr)f\le-f\T Rf\le-\lmin(R)\,\kappa^2\norm x^2<0 .
\]
(iii) Si $\D=\R^n$, la cota $V\ge\lmin(P)\kappa^2\norm x^2$ muestra que $V$ es radialmente no acotada \emph{automáticamente}, y se aplica el \cref{thm:ly-bk}.
\end{demostracion}

Tres observaciones. En $x=0$ la condición dice que $J(0)\T P+PJ(0)<0$, es decir, que la linealización es Hurwitz; el valor del método está en que la condición se pide \alerta{de manera uniforme en $\D$}, y eso da una región y, si $\D=\R^n$, un resultado global. La convexidad (basta que $\D$ sea estrellado respecto al origen) se usa solo en (i); sin ella, $V=\norm f_P^2$ puede anularse en otro equilibrio dentro de $\D$. Y las cotas de la prueba dan estabilidad \emph{exponencial}: cerca del origen $V\le\lmax(P)L^2\norm x^2$, con $L$ la constante de Lipschitz de $f$, y se aplica el \cref{thm:ly-exp}.

\begin{ejemplo}[label=ej:iv-k1]{El ejercicio del handout, resuelto}
Pruebe que el origen de $\dot x_1=-3x_1+x_2$, $\dot x_2=x_1-x_2-x_2^3$ es globalmente asintóticamente estable.
\solucion La jacobiana $J(x)=\bigl[\begin{smallmatrix}-3&1\\1&-1-3x_2^2\end{smallmatrix}\bigr]$ es simétrica, así que con $P=I$, $J\T P+PJ=2J$. Con $R=I$ hay que verificar $2J+I\le0$:
\[
  -(2J+I)=\begin{bmatrix}5&-2\\-2&1+6x_2^2\end{bmatrix},
\]
con menores principales $5>0$ y $5(1+6x_2^2)-4=1+30x_2^2>0$: es definida positiva para \emph{todo} $x\in\R^2$. Por el \cref{thm:iv-krasovskii} con $\D=\R^2$, el origen es globalmente asintóticamente estable, con
\[
  V(x)=\norm{f(x)}^2=(-3x_1+x_2)^2+(x_1-x_2-x_2^3)^2,\qquad \Vdot=2f\T Jf\le-\norm f^2 .
\]
Los niveles de $V$ no son elipses: el término cúbico los deforma (\cref{fig:iv-krasovskii}, izquierda). En el plano de velocidades $z=f(x)$, todas las soluciones se acercan al origen con $\norm{z(t)}$ decreciente (derecha).
\end{ejemplo}

\begin{figure}[ht]
\centering
\includegraphics[width=\linewidth]{cap08/f_krasovskii}
\caption{\Cref{ej:iv-k1}. Izquierda: niveles de $V=\norm{f(x)}^2$ y soluciones desde un círculo de radio $2.2$. Derecha: las mismas soluciones en el plano de velocidades $z=f(x)$; los círculos grises son niveles de $\norm z$.}
\label{fig:iv-krasovskii}
\end{figure}

\codigorepo{cap08_invarianza/C8_Ej12_Krasovskii.m}{k1}{La condición \eqref{eq:iv-krasovskii} en una malla y el plano de velocidades}

\begin{trampa}[label=tr:iv-PI]{``Krasovskii con $P=I$ falla, luego el origen no es asintóticamente estable''}
\plausible{$P=I$ es la elección natural; si no funciona, el sistema no cumple la condición.}
La condición es suficiente, y la elección de $P$ importa. Para el oscilador amortiguado $\dot x_1=x_2$, $\dot x_2=-x_1-x_2$, que es globalmente asintóticamente estable, con $P=I$ se tiene $J+J\T=\diag(0,-2)$: no existe $R>0$ con $J+J\T\le-R$. Con $P=\bigl[\begin{smallmatrix}1.5&0.5\\0.5&1\end{smallmatrix}\bigr]$, la solución de $A\T P+PA=-I$, la condición se cumple con $R=I$.
\correcto{Un fracaso con un $P$ particular no dice nada sobre la estabilidad. En sistemas lineales, el $P$ correcto sale de la ecuación de Lyapunov; en no lineales, se busca (a veces con desigualdades matriciales lineales).}
\end{trampa}

\codigorepo{cap08_invarianza/C8_Ej12_Krasovskii.m}{matrizP}{$P=I$ contra el $P$ de la ecuación de Lyapunov}

\begin{ejemplo}[label=ej:iv-k2]{$V=\norm f^2$ se anula en \emph{todos} los equilibrios}
Para $\dot x_1=-x_1+x_1^3$, $\dot x_2=-x_2$, con equilibrios en el origen (nodo estable) y en $(\pm1,0)$ (sillas) y región de atracción exacta $\{\abs{x_1}<1\}$, aplique Krasovskii con $P=I$.
\solucion $J+J\T=\diag(-2+6x_1^2,\ -2)$. En el dominio convexo $\D_r=\{\abs{x_1}<r\}$ con $r<1/\sqrt3$, la condición vale con $R=(2-6r^2)I>0$; en $\abs{x_1}=1/\sqrt3$ se pierde. En $\D_r$ el teorema da estabilidad asintótica con $V=(x_1-x_1^3)^2+x_2^2$. El mínimo de $V$ sobre $\abs{x_1}=r$ es $(r-r^3)^2$, que tiende a $4/27=0.148$ cuando $r\to1/\sqrt3$: para todo $c<4/27$, la componente de $\{V\le c\}$ que contiene al origen está en algún $\D_r$ y es un estimado válido.
\end{ejemplo}

\begin{trampa}[label=tr:iv-normaf]{``$V=\norm{f}^2$ es pequeña, luego el estado está cerca del origen''}
\plausible{$V$ es definida positiva y mide la distancia al equilibrio.}
$V(\pm1,0)=0$: $\{V\le c\}$ tiene \emph{tres} componentes, una alrededor de cada equilibrio (\cref{fig:iv-krasovskii-trampa}). El punto $(1.05,0)$ tiene $V=0.0116<0.12$ y escapa ($x_1=10$ en $t=1.18$). Fuera de $\D$, $V=\norm f^2$ no distingue un equilibrio estable de uno inestable: solo mide qué tan rápido se mueve el estado.
\correcto{$V=f\T Pf$ es definida positiva \emph{en $\D$}, donde el origen es el único equilibrio. El estimado es la componente del origen, y solo mientras esté dentro de $\D$.}
\end{trampa}

\begin{figure}[ht]
\centering
\includegraphics[width=0.78\linewidth]{cap08/f_krasovskii_trampa}
\caption{\Cref{ej:iv-k2}. Azul: la componente válida de $\{V\le c\}$ ($c=0.12$). Gris: las componentes alrededor de las sillas. Punteadas: $\abs{x_1}=1/\sqrt3$, el límite de los dominios $\D_r$. Coral: $\abs{x_1}=1$, la frontera de la región de atracción. La solución verde converge aunque está fuera del estimado; la coral escapa.}
\label{fig:iv-krasovskii-trampa}
\end{figure}

\codigorepo{cap08_invarianza/C8_Ej13_KrasovskiiTrampa.m}{k2}{$\norm f^2$ en las sillas y el estimado válido}

%% ---------------------------------------------------------------------
\section{La ecuación de Zubov: la región de atracción exacta}
\label{sec:iv-zubov}
%% ---------------------------------------------------------------------

El gradiente variable y Krasovskii dan \emph{estimados} conservadores de la región de atracción. Zubov cambia la pregunta: en lugar de buscar una $V$ cualquiera, impone la forma exacta de $\Vdot$ con una ecuación en derivadas parciales. Si la ecuación se resuelve con la condición de frontera adecuada, se obtiene la región de atracción \clave{exacta}.

\begin{teorema}[label=thm:iv-zubov]{Zubov \citep{Zub64}, \citep[\S8.2]{Kha02}}
Sea $x=0$ un equilibrio de $\dot x=f(x)$, con $f$ localmente Lipschitz, y $\D\subset\R^n$ un conjunto \clave{abierto} que contiene al origen. Suponga que existen $V:\D\to\R$ de clase $\Cont1$ y $h:\R^n\to\R$ continua tales que
\begin{enumerate}
  \item[(a)] $V(0)=0$ y $0<V(x)<1$ para $x\in\D\setminus\{0\}$;
  \item[(b)] $h(0)=0$ y $h(x)>0$ para $x\neq0$;
  \item[(c)] $V(x)\to1$ cuando $x\to\partial\D$, o cuando $\norm x\to\infty$ con $x\in\D$;
  \item[(d)] se cumple la \clave{ecuación de Zubov}
  \begin{equation}
    \jac Vx(x)\,f(x)=-h(x)\bigl(1-V(x)\bigr),\qquad x\in\D .
    \label{eq:iv-zubov}
  \end{equation}
\end{enumerate}
Entonces el origen es asintóticamente estable y \alerta{$\D$ es exactamente su región de atracción}. Si $\D=\R^n$, es globalmente asintóticamente estable.
\end{teorema}

\paragraph{La forma más fácil de leerlo.} Con $W=-\ln(1-V)$, que va de $\D$ en $[0,\infty)$, la ecuación de Zubov es
\begin{equation}
  \dot W=\frac{\Vdot}{1-V}=-h(x),\qquad W(x)\to\infty\ \text{ cuando }x\to\partial\D\ \text{ o }\ \norm x\to\infty .
  \label{eq:iv-W}
\end{equation}
$W$ es una función de Lyapunov que \alerta{explota exactamente en la frontera de la región de atracción}; la transformación $V=1-e^{-W}$ solo comprime $[0,\infty)$ en $[0,1)$.

\begin{demostracion}[Esbozo]
(i) \emph{Estabilidad asintótica}: $V$ es definida positiva en $\D$ y $\Vdot=-h(1-V)<0$ en $\D\setminus\{0\}$. (ii) \emph{$\D\subset\RA$}: si $x_0\in\D$, $W(x(t))\le W(x_0)<\infty$ mientras la solución esté en $\D$; como $W\to\infty$ en $\partial\D$ y en el infinito, la solución no puede alcanzar la frontera ni escapar, $\{W\le W(x_0)\}$ es compacto y $\dot W=-h<0$ fuera del origen: $x(t)\to0$. (iii) \emph{$\RA\subset\D$}: si $x_0\notin\D$ y $x(t)\to0$, la solución entra a $\D$ en algún $t_1$; yendo hacia atrás desde $t_1$, toca $\partial\D$ en algún $t_0<t_1$. Pero en $(t_0,t_1]$, $W(x(t))=W(x(t_1))+\int_t^{t_1}h(x(s))\dif s$ permanece acotada, mientras que $W\to\infty$ al acercarse a $\partial\D$: contradicción.
\end{demostracion}

\paragraph{Cómo se construye $V$.} Integrando $\dot W=-h$ a lo largo de la solución que sale de $x_0$ (método de las características),
\begin{equation}
  W(x_0)=\int_0^\infty h\bigl(\flujo t{x_0}\bigr)\dif t,\qquad V(x_0)=1-e^{-W(x_0)} .
  \label{eq:iv-caracteristicas}
\end{equation}
$W$ es finita exactamente cuando la solución converge al origen lo bastante rápido, es decir, en la región de atracción. Cerca del origen, con $h=x\T Qx$ y $A=J(0)$, se tiene $W(x)\approx x\T Px$ con $A\T P+PA=-Q$: \alerta{Zubov contiene a la ecuación de Lyapunov como su aproximación lineal}. Y $W$ con $h=\norm x^2$ es la ``energía futura'' del converso de la \cref{sec:av-conversos}.

\begin{ejemplo}[label=ej:iv-z1]{Zubov en la recta}
Para $\dot x=-x+x^3$ (equilibrios $0$ estable y $\pm1$ inestables, $\RA=(-1,1)$), resuelva la ecuación de Zubov con $h=2x^2$.
\solucion $V'(x)(-x+x^3)=-2x^2(1-V)$ es separable:
\[
  \frac{\dif V}{1-V}=\frac{2x^2}{x-x^3}\dif x=\frac{2x}{1-x^2}\dif x\implica-\ln(1-V)=-\ln(1-x^2) .
\]
Luego $V=x^2$ y $W=-\ln(1-x^2)$. $V<1$ exactamente en $(-1,1)$ y $V\to1$ en $\pm1$: el teorema entrega la región de atracción exacta (\cref{fig:iv-zubov1d}).
\end{ejemplo}

\begin{figure}[ht]
\centering
\includegraphics[width=0.9\linewidth]{cap08/f_zubov1d}
\caption{\Cref{ej:iv-z1}. Izquierda: $f(x)=-x+x^3$. Derecha: la solución de Zubov $V=x^2$ y $W=-\ln(1-V)$, que explota en la frontera $\pm1$ de la región de atracción.}
\label{fig:iv-zubov1d}
\end{figure}

\begin{ejemplo}[label=ej:iv-z2]{Una región de atracción no acotada}
Encuentre la región de atracción del origen de $\dot x_1=-x_1+2x_1^2x_2$, $\dot x_2=-x_2$ y certifíquela con Zubov.
\solucion \emph{Identificar la región.} La cantidad $w=x_1x_2$ cumple
\[
  \dot w=\dot x_1x_2+x_1\dot x_2=-2w+2w^2=-2w(1-w) .
\]
Si $w(0)<1$, $w\to0$ y luego $x\to0$; la curva $w=1$ es invariante, y si $w(0)>1$, $w$ explota. La candidata es $\D=\{x_1x_2<1\}$, acotada por la hipérbola.
\emph{Solución de Zubov.} Proponga $W=\tfrac12x_2^2+\dfrac{x_1^2}{2(1-w)}$. Con $\dot x_1=-x_1(1-2w)$,
\[
  \jac W{x_1}\dot x_1=\Bigl(\frac{x_1}{1-w}+\frac{x_1w}{2(1-w)^2}\Bigr)\bigl(-x_1(1-2w)\bigr),\qquad
  \jac W{x_2}\dot x_2=-x_2^2-\frac{x_1^2w}{2(1-w)^2} .
\]
Al sumar, los términos en $x_1^2$ se agrupan como
\[
  -x_1^2\Bigl[\frac{1-2w}{1-w}+\frac{w(1-2w)+w}{2(1-w)^2}\Bigr]=-x_1^2\Bigl[\frac{1-2w}{1-w}+\frac w{1-w}\Bigr]=-x_1^2 ,
\]
y $\dot W=-(x_1^2+x_2^2)=-h(x)$. $W\to\infty$ cuando $x_1x_2\to1^-$ y cuando $\norm x\to\infty$ dentro de $\D$: con $V=1-e^{-W}$ se cumplen todas las hipótesis del \cref{thm:iv-zubov}, y $\RA=\{x_1x_2<1\}$. Desde $(1.9,0.5)$ ($w=0.95$) la solución converge; desde $(2.1,0.5)$ ($w=1.05$) escapa en tiempo finito (\cref{fig:iv-zubov-hiperbola}).
\end{ejemplo}

\begin{figure}[ht]
\centering
\includegraphics[width=0.68\linewidth]{cap08/f_zubov_hiperbola}
\caption{\Cref{ej:iv-z2}. Niveles $V=0.25$, $0.5$, $0.75$, $0.9$ y $0.99$ (azul); en gris, el exterior de la región de atracción; en coral, $x_1x_2=1$. Las soluciones negras empiezan dentro; la coral punteada, desde $(2.1,0.5)$, fuera.}
\label{fig:iv-zubov-hiperbola}
\end{figure}

\codigorepo{cap08_invarianza/C8_Ej14_Zubov.m}{z2}{El residuo de la ecuación de Zubov en una malla y los niveles $V\to1$}

\begin{ejemplo}[label=ej:iv-z3]{Zubov numérico: van der Pol en tiempo invertido}
Para \eqref{eq:ly-vdpinv}, $\dot x_1=-x_2$, $\dot x_2=x_1+(x_1^2-1)x_2$, la frontera de la región de atracción es el ciclo límite inestable. Calcule la función de Zubov numéricamente y compárela con el mejor estimado cuadrático del \cref{ej:ly-vdp}.
\solucion No hay solución cerrada de \eqref{eq:iv-zubov}, pero con $h=\norm x^2$ se calcula \eqref{eq:iv-caracteristicas} simulando desde cada punto de una malla, con $W=\infty$ si la solución no converge. Cerca del origen, $W(0.1,0)=0.0150$, que coincide con $x\T Px$ para la $P=\bigl[\begin{smallmatrix}1.5&-0.5\\-0.5&1\end{smallmatrix}\bigr]$ del capítulo~7; en $(1,0)$, $W=1.60$ contra $x\T Px=1.5$; en $(1.9,0)$, cerca del ciclo, $W=11.2$; en $(2.1,0)$, fuera, $W=\infty$. La región de atracción (el interior del ciclo) tiene área $13.7$; la mejor elipse $x\T Px\le2.30$ tiene área $6.46$, el 47\,\%. Los niveles de Zubov llenan toda la región hasta el ciclo (\cref{fig:iv-zubov-vdp}).
\end{ejemplo}

\begin{figure}[ht]
\centering
\includegraphics[width=\linewidth]{cap08/f_zubov_vdp}
\caption{\Cref{ej:iv-z3}. Izquierda: niveles $V=0.5$, $0.9$, $0.99$ y $0.999$ de la función de Zubov calculada en una malla (azul), el ciclo límite (coral) y el mejor estimado cuadrático (verde). Derecha: $V(x)$ en la malla; en blanco, los puntos con $W=\infty$.}
\label{fig:iv-zubov-vdp}
\end{figure}

\codigorepo{cap08_invarianza/C8_Ej15_ZubovVanDerPol.m}{zubov}{$W(x_0)=\int_0^\infty\norm{\flujo t{x_0}}^2\dif t$ en una malla}

\video{V8_4_zubov_vdp}{Los niveles de Zubov llenan la región hasta el ciclo}

\begin{trampa}[label=tr:iv-frontera]{``$V$ resuelve la ecuación de Zubov en $\D$, luego $\D$ es la región de atracción''}
\plausible{La ecuación \eqref{eq:iv-zubov} es la condición del teorema, y $V$ la cumple.}
Falta la condición de frontera (c): $V\to1$ en \emph{toda} $\partial\D$ y en el infinito. Una solución de la EDP que no tiende a $1$ en parte de la frontera solo certifica que esa parte \emph{no} es frontera de $\RA$. Por ejemplo, en el \cref{ej:iv-z2}, $V=1-e^{-W}$ resuelve la ecuación en cualquier $\D'=\{x_1x_2<1,\ \norm x<5\}$, pero en el círculo $\norm x=5$ no tiende a $1$, y $\D'$ no es la región de atracción.
\correcto{La EDP fija $V$ en el interior; la condición de frontera es la que identifica la región.}
\end{trampa}

\begin{profundiza}[Qué da Zubov y qué no]
\emph{Sí}: caracteriza la región de atracción exacta, y diferentes $h$ producen diferentes $V$ pero el mismo conjunto $\{V<1\}$. \emph{No}: rara vez se resuelve en forma cerrada; calcular $W$ numéricamente equivale a simular desde cada punto, así que su valor es conceptual y como certificado, más que como atajo computacional. Zubov mismo propuso desarrollos en series; hoy se buscan funciones polinomiales con \emph{sumas de cuadrados} (SOS), que convierten ``$V\ge0$'' y ``$-\Vdot\ge0$'' en problemas de programación semidefinida~\citep{Par00}. Es la versión moderna de las dos preguntas de este capítulo: encontrar $V$ y certificar una región.
\end{profundiza}

%% ---------------------------------------------------------------------
\section{Para llevar}
\label{sec:iv-cierre}
%% ---------------------------------------------------------------------

\begin{parallevar}
\begin{enumerate}
  \item \textbf{Invarianza:} un conjunto invariante está hecho de soluciones completas; un conjunto de nivel con $\Vdot\le0$ es solo positivamente invariante, y el que sirve es la \emph{componente conexa} que contiene la condición inicial.
  \item \textbf{LaSalle:} con $\Omega$ compacto y positivamente invariante y $\Vdot\le0$ en $\Omega$, toda solución tiende a $M$, el mayor conjunto invariante en $E=\{\Vdot=0\}$. La prueba es la cadena $\omegalim{x_0}\subset M\subset E\subset\Omega$. $V$ no tiene que ser definida positiva.
  \item \textbf{$M$ se calcula con la prueba de permanencia:} derivar las restricciones que definen $E$ hasta que obliguen a la solución a algo conocido. Para sistemas lineales es la prueba de observabilidad.
  \item \textbf{LaSalle da convergencia a un conjunto}, no a un punto, y no da estabilidad. Los corolarios de Barbashin y Krasovskii la recuperan pidiendo $V$ definida positiva (y radialmente no acotada para el resultado global).
  \item \textbf{Construir $V$:} el gradiente variable fija $\nabla V$ y deja parámetros para optimizar el estimado; Krasovskii fija $V=f\T Pf$ y da un resultado global con una sola condición matricial; Zubov fija $\Vdot$ y da la región de atracción exacta, casi siempre solo numéricamente.
\end{enumerate}
\end{parallevar}

\begin{center}
\small
\begin{tabular}{@{}p{2.4cm}p{4.4cm}p{5.0cm}@{}}
\toprule
& Fortaleza & Debilidad \\
\midrule
Gradiente variable & sistemático; parámetros libres para mejorar el estimado & la forma de $g$ hay que adivinarla; el álgebra crece rápido con $n$ \\
Krasovskii & una sola condición matricial; estabilidad global de un golpe & muy exigente (uniforme en $\D$); $V=\norm f_P^2$ no distingue equilibrios \\
Zubov & región de atracción exacta & EDP sin solución cerrada en general \\
\bottomrule
\end{tabular}
\end{center}

\begin{trampa}[title={\faBug\hspace{0.45em}Las trampas del capítulo, en una caja}]{}
\begin{itemize}
  \item $\{V\le c\}$ no es lo mismo que su componente del origen (\cref{tr:iv-subnivel,tr:iv-componente}).
  \item $E$ no es $M$ (\cref{tr:iv-EM}); revisar solo los equilibrios de $E$ no basta (\cref{tr:iv-equilibrios}); pasar por $S$ no es permanecer en $S$ (\cref{tr:iv-pasar}).
  \item Global sin radialidad es falso (\cref{tr:ly-radial}); convergencia a $M$ no es convergencia a un punto (\cref{tr:iv-punto}) ni estabilidad (\cref{tr:iv-estabilidad}).
  \item LaSalle es para sistemas invariantes en el tiempo; con $\dot x=f(t,x)$ se usan Barbalat o LaSalle--Yoshizawa.
  \item Hay que exhibir $\Omega$ compacto y positivamente invariante antes de concluir.
  \item Un estimado no es la región (\cref{tr:iv-estimado}); la elección limpia no es la óptima (\cref{tr:iv-limpia}); un $g$ no simétrico no es un gradiente (\cref{tr:iv-simetria}).
  \item Krasovskii con un $P$ que falla no dice nada (\cref{tr:iv-PI}); $\norm f^2$ se anula en todos los equilibrios (\cref{tr:iv-normaf}); Zubov sin condición de frontera no certifica nada (\cref{tr:iv-frontera}).
\end{itemize}
\end{trampa}

\begin{profundiza}[Correcciones respecto a la versión anterior (LaSalle\_Invarianza.tex, Construccion\_Lyapunov.tex y Handout7)]
(i) El teorema de LaSalle pedía ``$\Omega$ compacto e \emph{invariante}''; lo correcto, y lo que se usa con los conjuntos de nivel, es \emph{positivamente} invariante. (ii) El corolario global omitía que $V$ debe ser radialmente no acotada. (iii) ``La única trayectoria que está en $S$'' se cambió por ``ninguna solución distinta de $x\equiv0$ puede permanecer idénticamente en $S$''. (iv) En el gradiente variable, la condición de simetría del ejemplo decía $\partial g_1/\partial x_1=\partial g_2/\partial x_1$; es $\partial g_1/\partial x_2=\partial g_2/\partial x_1$. (v) Krasovskii: ``la solución $x(t)=0$ es única'' se cambió por ``el origen es el único equilibrio en $\D$'', con $\D$ convexo, y se agregó la prueba de que $f\T Pf$ es definida positiva y, si $\D=\R^n$, radialmente no acotada. (vi) Zubov: $V:\D\to\R$ y $h:\R^n\to\R$ (no $\R^n$); $\D$ abierto (no ``acotado'', que era incompatible con la versión global), y la conclusión es que $\D$ \emph{es} la región de atracción, no solo que está contenida en ella. (vii) El ejemplo de la $V$ no radialmente no acotada, que estaba repetido, ahora se cita del capítulo~7, y el ejemplo del gradiente variable para el péndulo se presenta como la deducción de la $V$ del \cref{ej:ly-cruzado} (con $c^*=2.308$ para $b=0.5$). (viii) El mejor $a_{11}$ del \cref{ej:iv-g1} es $1.48$ (área $37.4$), no ``cerca de $1.5$''. (ix) La trampa de sistemas variantes en el tiempo remitía a ``más adelante en este capítulo (lámina 147)''; LaSalle--Yoshizawa se usa en el \cref{cap:variantes} y Barbalat en el \cref{cap:avanzada}. (x) Notación: invariante en el tiempo en lugar de autónomo; $\Nivel c$ es siempre una componente conexa; $\flujo t{x_0}$ para las soluciones; Lotka--Volterra con la notación del \cref{cap:orbitas}; el ejemplo de observabilidad usa las dos masas del \cref{cap:lyapunov}.
\end{profundiza}

%% ---------------------------------------------------------------------
\section{Problemas}
\label{sec:iv-problemas}
%% ---------------------------------------------------------------------

Los problemas marcados con \etiquetaM{} requieren \matlab/\simulink; los marcados con \etiquetaEstrella{} son de profundización. Cada uno tiene una trampa: búsquela antes de dar la respuesta.

\subsection*{Invarianza y LaSalle}

\begin{problema}[sistemas de Liénard]\label{prob:iv-lienard}
Para $\dot x_1=x_2$, $\dot x_2=-g(x_1)-k(x_2)$ con $g,k$ continuas, $sg(s)>0$ y $sk(s)>0$ para $s\neq0$: (a) con $V=\int_0^{x_1}g(s)\dif s+\tfrac12x_2^2$, calcule $\Vdot$; (b) pruebe que el origen es asintóticamente estable; (c) ¿qué condición sobre $g$ da estabilidad asintótica global? (d) ¿Qué falla con $g(s)=s/(1+s^2)^2$?
\end{problema}

\begin{problema}[un sistema gradiente]\label{prob:iv-gradiente}
Sea $\dot x=-\nabla U(x)$ con $U=\tfrac14(x_1^2-1)^2+\tfrac12x_2^2$. (a) Con $V=U$, muestre que toda solución converge al conjunto $\{(0,0),(\pm1,0)\}$; (b) ¿converge a un punto? (c) ¿Cuáles equilibrios son estables? (d) ¿Converge toda solución a un mínimo de $U$?
\end{problema}

\begin{problema}[el péndulo en el cilindro]\label{prob:iv-cilindro}
Formalice la observación de la \cref{sec:iv-lasalle}: muestre que toda solución del péndulo con fricción converge a $(0,0)$ o a $(\pi,0)$, módulo $2\pi$. ¿Por qué no se puede trabajar con $\{V\le c\}\subset\R^2$?
\end{problema}

\begin{problema}[un circuito con zona muerta]\label{prob:iv-rlc}
Considere $L\,\dot\imath=-v_C-\psi(i)$, $C\,\dot v_C=i$, con $\psi(i)=0$ si $\abs i\le\delta$ y $\psi(i)=R\,(i-\delta\sgn i)$ si $\abs i>\delta$. (a) Con $V=\tfrac L2i^2+\tfrac C2v_C^2$, muestre que $\Vdot=-i\,\psi(i)\le0$; (b) calcule $E$ y $M$, y muestre que $M=\{V\le L\delta^2/2\}$; (c) ¿qué se puede concluir? ¿Se descarga completamente el circuito?
\end{problema}

\begin{problema}[el amortiguador en la otra masa]\label{prob:iv-masas}
\etiquetaM{} Para las dos masas del \cref{ej:iv-observable}: (a) con $\kappa=0$, calcule $\ker\mathcal O$ y verifique que coincide con $M$; (b) mueva el amortiguador a la masa~2 y repita; (c) con $\kappa=0.4$, ¿cambia la estabilidad si el amortiguador está en una u otra masa?
\end{problema}

\begin{problema}[el límite del parámetro adaptativo]\label{prob:iv-adaptativo}
\etiquetaM{} Ejecute \cmd{C8\_Ej5\_Adaptativo.m} con otras condiciones iniciales y compare $k(T)$ con \eqref{eq:iv-kinf}. (a) ¿Qué pasa con la fórmula si $y_0=0$? (b) Construya el modelo con \cmd{crearModelo\_AdaptativoLaSalle.m} y verifique que $k_\infty$ no depende de la constante $b$ del análisis.
\end{problema}

\subsection*{Construcción de funciones de Lyapunov}

\begin{problema}[gradiente variable y Zubov sobre el mismo sistema]\label{prob:iv-slotine}
Para $\dot x_1=-2x_1$, $\dot x_2=-2x_2+2x_1x_2^2$~\citep{SlLi91}: (a) con $g=(x_1,x_2)$, muestre que $\Vdot=-2x_1^2-2x_2^2(1-x_1x_2)$ y que el estimado es $x_1^2+x_2^2<2$; (b) con $w=x_1x_2$, muestre que la región de atracción exacta es $\{x_1x_2<2\}$; (c) encuentre una solución de Zubov inspirada en el \cref{ej:iv-z2}. ¿Es globalmente asintóticamente estable el origen, ya que $V$ es radialmente no acotada?
\end{problema}

\begin{problema}[el mejor término cruzado, otra vez]\label{prob:iv-pendulo}
En el \cref{ej:iv-g2}, verifique que $c^*(k)$ es máximo en $k=b/2$ y compare el estimado con el que da LaSalle con la energía ($c<2$). ¿Cuál es más grande? ¿Cuál exige un argumento adicional?
\end{problema}

\begin{problema}[Krasovskii global]\label{prob:iv-krasovskii}
Muestre con $P=I$ que el origen de $\dot x_1=-6x_1+2x_2$, $\dot x_2=2x_1-6x_2-2x_2^3$ es globalmente asintóticamente estable. ¿Hace falta verificar que $V$ es radialmente no acotada?
\end{problema}

\begin{problema}[escoger $P$]\label{prob:iv-escogerP}
Para $\dot x_1=x_2$, $\dot x_2=-x_1-x_2$: (a) muestre que Krasovskii con $P=I$ falla; (b) encuentre un $P$ con el que funcione; (c) escriba la $V$ que resulta.
\end{problema}

\begin{problema}[Zubov en la recta, con una región no acotada]\label{prob:iv-zubov1d}
Para $\dot x=-x+x^2$ con $h=2x^2$, encuentre $W$ y $V$. Verifique que $V\to1$ cuando $x\to1^-$ y cuando $x\to-\infty$, y concluya que la región de atracción es $(-\infty,1)$.
\end{problema}

\begin{problema}[el mismo conjunto con otra $V$]\label{prob:iv-zubovh}
\etiquetaM{} Modifique \cmd{C8\_Ej15\_ZubovVanDerPol.m} con $h=2x_1^2+x_2^2$. Compare los niveles de $V$ con los de $h=\norm x^2$ y verifique que $\{V<1\}$ no cambia.
\end{problema}

%% ---------------------------------------------------------------------
\section{Soluciones seleccionadas}
%% ---------------------------------------------------------------------

\paragraph{Problema~\ref{prob:iv-lienard}}
(a) $\Vdot=g(x_1)x_2+x_2(-g(x_1)-k(x_2))=-x_2k(x_2)\le0$. (b) $V$ es definida positiva cerca del origen y $S=\{x_2=0\}$; $x_2\equiv0$ obliga $g(x_1)\equiv0$, es decir $x_1\equiv0$: por el \cref{cor:iv-as}, asintóticamente estable. (c) Si $\int_0^sg\to\infty$ cuando $\abs s\to\infty$, $V$ es radialmente no acotada y el \cref{cor:iv-gas} da el resultado global. (d) $\int_0^sg=\tfrac12\bigl(1-1/(1+s^2)\bigr)<\tfrac12$: $V$ no es radialmente no acotada y la conclusión es local, con $\{V<\tfrac12\}$ como estimado. La trampa es concluir global solo porque $\Vdot\le0$ en todo $\R^2$.

\paragraph{Problema~\ref{prob:iv-gradiente}}
(a) $\dot U=-\norm{\nabla U}^2\le0$; $U$ es radialmente no acotada, así que los $\{U\le c\}$ son compactos; $E=\{\nabla U=0\}$ son los tres equilibrios, y $M=E$. (b) Sí: los equilibrios son aislados y $\omegalim{x_0}$ es conexo (\cref{lem:oc-omega}), así que es uno solo. (c) $(\pm1,0)$ son mínimos de $U$ y estables; $(0,0)$ es una silla (hessiana $\diag(-1,1)$), inestable. (d) No: las soluciones que nacen en el eje $x_1=0$, la variedad estable de la silla, convergen a $(0,0)$.

\paragraph{Problema~\ref{prob:iv-cilindro}}
En $\mathbb S^1\times\R$, $V=a(1-\cos x_1)+\tfrac12x_2^2$ es continua y $\{V\le c\}$ es compacto para todo $c$; $E=\{x_2=0\}$ y la permanencia da $\sin x_1\equiv0$: $M=\{(0,0),(\pi,0)\}$. Como los dos puntos son aislados y $\omegalim{x_0}$ es conexo, cada solución converge a uno de ellos. En $\R^2$, $\{V\le c\}$ no es acotado para ningún $c>0$ y no hay compacto que usar.

\paragraph{Problema~\ref{prob:iv-rlc}}
(a) $\Vdot=Li\dot\imath+Cv_C\dot v_C=-iv_C-i\psi(i)+v_Ci=-i\psi(i)\le0$. (b) $E=\{\abs i\le\delta\}$. En $E$, $\psi=0$ y el circuito es un $LC$ sin pérdidas: $V$ es constante y la corriente oscila con amplitud $\sqrt{2V/L}$. Una solución se queda en $E$ si y solo si esa amplitud es $\le\delta$, es decir $V\le L\delta^2/2$. (c) Las soluciones tienden a $M$: la energía converge a algún valor $\le L\delta^2/2$, y el circuito sigue oscilando. No se descarga: la zona muerta deja una oscilación residual.

\paragraph{Problema~\ref{prob:iv-masas}}
(a) $\ker\mathcal O=\{q_1=v_1=0\}$, que es invariante (la masa~2 desacoplada) y es $M$. (b) Por simetría, $\ker\mathcal O=\{q_2=v_2=0\}$: la masa~1 oscila. (c) No: con $\kappa\neq0$ el par es observable en los dos casos y $A$ es Hurwitz.

\paragraph{Problema~\ref{prob:iv-adaptativo}}
(a) Si $y_0=0$, la solución es el equilibrio $(0,k_0)$ y $k_\infty=k_0$; la fórmula da $a+\abs{k_0-a}$, que coincide solo si $k_0\ge a$. (b) $b$ no aparece en el modelo: el lazo solo tiene $a$, $\gamma$ y la condición inicial.

\paragraph{Problema~\ref{prob:iv-slotine}}
(a) $\Vdot=-2x_1^2+x_2(-2x_2+2x_1x_2^2)=-2x_1^2-2x_2^2(1-x_1x_2)$; como $x_1x_2\le\tfrac12\norm x^2$, $\Vdot<0$ si $\norm x^2<2$, y el mayor nivel dentro es $\norm x^2<2$. (b) $\dot w=-2x_1x_2+x_1(-2x_2+2x_1x_2^2)=-2w(2-w)$: si $w(0)<2$, $w\to0$ y $x\to0$; si $w(0)>2$, $w$ explota. (c) $W=\tfrac12x_1^2+\dfrac{x_2^2}{2(1-w/2)}$ cumple $\dot W=-2\norm x^2$ y $W\to\infty$ cuando $w\to2^-$. (d) No es global: que $V$ sea radialmente no acotada no basta si $\Vdot$ no es negativa en todo $\R^2$ (desde $(3,1)$, $w=3$, la solución escapa).

\paragraph{Problema~\ref{prob:iv-pendulo}}
$\frac{\dif c^*}{\dif k}=\frac{\pi^2}2(b-2k)=0$ en $k=b/2$, con $c^*=2+\pi^2b^2/8$ ($2.308$ con $b=0.5$). LaSalle con la energía certifica los niveles $c<2$; la $V$ inclinada, niveles de otra función, con un área mayor. LaSalle necesita la prueba de permanencia; la $V$ estricta no.

\paragraph{Problema~\ref{prob:iv-krasovskii}}
$J=\bigl[\begin{smallmatrix}-6&2\\2&-6-6x_2^2\end{smallmatrix}\bigr]$ es simétrica y $-(2J+I)=\bigl[\begin{smallmatrix}11&-4\\-4&11+12x_2^2\end{smallmatrix}\bigr]$, con determinante $105+132x_2^2>0$: definida positiva en todo $\R^2$. No hace falta verificar la radialidad: el \cref{thm:iv-krasovskii} la da automáticamente ($V\ge\lmin(P)\kappa^2\norm x^2$).

\paragraph{Problema~\ref{prob:iv-escogerP}}
(a) $J+J\T=\diag(0,-2)$, que no es $\le-R$ para ningún $R>0$. (b) $P=\bigl[\begin{smallmatrix}1.5&0.5\\0.5&1\end{smallmatrix}\bigr]$, solución de $A\T P+PA=-I$. (c) $V=x\T A\T PA\,x=x_1^2+x_1x_2+1.5\,x_2^2$.

\paragraph{Problema~\ref{prob:iv-zubov1d}}
$\frac{\dif W}{\dif x}(-x+x^2)=-2x^2$ da $W'=\frac{2x}{1-x}$ y $W=-2x-2\ln(1-x)$, con $V=1-(1-x)^2e^{2x}$. $W'$ tiene el signo de $x$ en $x<1$, así que $W\ge W(0)=0$. Cuando $x\to1^-$, $\ln(1-x)\to-\infty$ y $V\to1$; cuando $x\to-\infty$, $(1-x)^2e^{2x}\to0$ y $V\to1$ (con $x=-5$, $V=1-36e^{-10}=0.998$). $\RA=(-\infty,1)$, no acotada.

\paragraph{Problema~\ref{prob:iv-zubovh}}
Los niveles intermedios cambian de forma (pesan más la dirección $x_1$), pero $W=\infty$ exactamente en los mismos puntos de la malla: $\{V<1\}$ es el interior del ciclo con cualquier $h$ definida positiva.

\respuestaspausas

